\documentclass[12pt]{article}

% === CÀI ĐẶT TRANG ===
\usepackage[margin=2cm]{geometry}
\usepackage[utf8]{inputenc}
\usepackage[T5]{fontenc}
\usepackage[vietnamese]{babel}

% === TOÁN HỌC ===
\usepackage{amsmath, amssymb, amsthm}

% === HÌNH VẼ ===
\usepackage{graphicx}
\usepackage{float}
\usepackage{tikz}
\usepackage{pgfplots}
\pgfplotsset{compat=1.18}

% === ĐỊNH DẠNG ===
\usepackage{enumitem}
\usepackage{xcolor}
\usepackage[most]{tcolorbox}
\usepackage{multicol}
\usepackage{booktabs}
\usepackage{array}
\usepackage{fancyhdr}
\usepackage{tocloft}
\usepackage[hidelinks]{hyperref}

% === LỆNH TỰ ĐỊNH NGHĨA ===
\newcommand{\otraloi}{\underline{\hspace{5cm}}}
\setlength{\parindent}{0pt} % Tắt lùi đầu dòng toàn giáo án

% === TCOLORBOX STYLES ===
\tcbuselibrary{breakable, skins, fitting}

% Ví dụ và Lời giải
\newtcolorbox{vidu}[1][1]{
	colback=blue!3!white, colframe=blue!55!black,
	fonttitle=\bfseries, title={Ví dụ #1},
	breakable, enhanced
}

\newtcolorbox{loigiai}{
	colback=gray!4!white, colframe=gray!40!black,
	fonttitle=\bfseries, title={Lời giải chi tiết},
	breakable, enhanced
}

% Hộp phương pháp giải nhanh
\newtcolorbox{phuongphap}[1][Mẹo giải nhanh Casio / Quy tắc nhận diện]{
	colback=orange!5!white, colframe=orange!75!black,
	fonttitle=\bfseries, title={#1},
	breakable, enhanced
}

% --- KHUNG BÀI TẬP TỰ LUẬN PHÂN MÀU ---
% Bài tập xanh – Bắt buộc
\newtcolorbox{btxanh}[1][]{
	colback=blue!4!white, colframe=blue!60!black,
	fonttitle=\bfseries\color{white}, coltitle=white,
	attach boxed title to top left={yshift=-2mm, xshift=4mm},
	boxed title style={colback=blue!60!black, rounded corners}, breakable, enhanced, #1
}

% Bài tập vàng – Bổ sung
\newtcolorbox{btvang}[1][]{
	colback=yellow!10!white, colframe=orange!70!black,
	fonttitle=\bfseries, coltitle=orange!80!black,
	attach boxed title to top left={yshift=-2mm, xshift=4mm},
	boxed title style={colback=yellow!40!white, colframe=orange!70!black, rounded corners}, breakable, enhanced, #1
}

% Bài tập đỏ – Gia cố
\newtcolorbox{btdo}[1][]{
	colback=red!4!white, colframe=red!60!black,
	fonttitle=\bfseries\color{white}, coltitle=white,
	attach boxed title to top left={yshift=-2mm, xshift=4mm},
	boxed title style={colback=red!60!black, rounded corners}, breakable, enhanced, #1
}

% === HEADER/FOOTER ===
\pagestyle{fancy}
\fancyhf{}
\fancyhead[L]{\small Giáo án Toán 12}
\fancyhead[C]{\small \textbf{Chủ đề: Đường tiệm cận của đồ thị hàm số}}
\fancyhead[R]{\small Lớp: 12}
\fancyfoot[C]{\thepage}
\renewcommand{\headrulewidth}{0.4pt}

% ================================================
\begin{document}
	
	\begin{center}
		{\LARGE \textbf{ĐƯỜNG TIỆM CẬN CỦA ĐỒ THỊ HÀM SỐ}}\\[6pt]
		{\large Môn: Toán 12 \quad Thời lượng: 1 tiết}
	\end{center}
	\vspace{4mm}
	\hrule
	\vspace{4mm}
	
	% ===========================
	% PHẦN 1: MỤC TIÊU BÀI HỌC
	% ===========================
	\section*{Mục tiêu bài học}
	\addcontentsline{toc}{section}{Mục tiêu bài học}
	
	\textbf{1. Kiến thức:}
	\begin{itemize}
		\item Nắm vững định nghĩa và ý nghĩa hình học của đường tiệm cận đứng (TCĐ), đường tiệm cận ngang (TCN) và đường tiệm cận xiên (TCX).
		\item Biết cách xác định nhanh các đường tiệm cận của đồ thị hàm số phân thức bậc nhất/bậc nhất, bậc hai/bậc nhất bằng cả tự luận và máy tính Casio.
	\end{itemize}
	
	\textbf{2. Kĩ năng:}
	\begin{itemize}
		\item Kĩ năng tính giới hạn vô cực ($\lim_{x \to x_0^{\pm}} f(x) = \pm\infty$ và $\lim_{x \to \pm\infty} f(x) = y_0$) để xác định tiệm cận.
		\item Vận dụng thành thạo kĩ năng chia đa thức và bấm máy tính để tìm nhanh phương trình tiệm cận xiên $y = ax + b$.
	\end{itemize}
	
	\textbf{3. Thái độ:}
	\begin{itemize}
		\item Tư duy hình học trực quan tốt, cẩn thận trong việc đối chiếu nghiệm của tử và mẫu để tránh bẫy triệt tiêu nghiệm.
	\end{itemize}
	
	\newpage
	
	% =====================
	% PHẦN 2: MỤC LỤC
	% =====================
	\setcounter{tocdepth}{2}
	\tableofcontents
	\newpage
	
	% =====================
	% PHẦN 3: LÝ THUYẾT
	% =====================
	\section{Lý thuyết và Phương pháp giải nhanh}
	
	\subsection{Đường tiệm cận ngang (TCN)}
	
	Đường tiệm cận ngang của đồ thị hàm số biểu diễn giới hạn của hàm số khi biến số $x$ tiến ra vô cực ($+\infty$ hoặc $-\infty$).
	
	\begin{phuongphap}[Định nghĩa toán học]
		Đường thẳng $y = y_0$ được gọi là đường \textbf{tiệm cận ngang} của đồ thị hàm số $y = f(x)$ nếu thỏa mãn ít nhất một trong các điều kiện sau:
		\[ \lim_{x \to +\infty} f(x) = y_0 \quad \text{hoặc} \quad \lim_{x \to -\infty} f(x) = y_0 \]
	\end{phuongphap}
	
	% === HÌNH VẼ MINH HỌA LÝ THUYẾT TỔNG QUÁT ===
	\begin{figure}[H]
		\centering
		\begin{tikzpicture}
			\begin{axis}[
				axis lines = middle,
				xlabel = {$x$},
				ylabel = {$y$},
				xmin = -1, xmax = 8,
				ymin = -1, ymax = 5,
				xtick = \empty,
				ytick = {3},
				yticklabels = {$y_0$},
				axis line style={->},
				width = 10cm, height = 6.5cm,
				clip = false
				]
				% Đường tiệm cận ngang y = y0 (với y0 = 3)
				\addplot [
				domain=-1:8, 
				dashed, 
				color=black!60, 
				thick
				] {3};
				
				% Đồ thị hàm số y = f(x) tiếp cận y0 từ phía dưới (hàm mũ giảm dần khoảng cách)
				\addplot [
				domain=0.2:8, 
				samples=100, 
				color=blue,
				thick
				] {3 - 2.8*exp(-0.5*x)};
				
				% Nhãn cho đồ thị và đường tiệm cận
				\node[blue, right] at (axis cs:3.5, 2.2) {$y = f(x)$};
				\node[black!80, above] at (axis cs:6, 3) {$y = y_0$ (Tiệm cận ngang)};
				
				% Mũi tên chỉ hướng biến x tiến ra dương vô cực
				\draw[->, thick, black!40] (axis cs:2, 0.4) -- (axis cs:6, 0.4) node[midway, below] {$x \to +\infty$};
			\end{axis}
		\end{tikzpicture}
		\caption{Minh họa hình học của đường tiệm cận ngang $y = y_0$ khi $x$ tiến ra $+\infty$}
	\end{figure}
	
	\begin{phuongphap}[Quy tắc so sánh bậc của Tử số $P(x)$ và Mẫu số $Q(x)$]
		Xét hàm số phân thức $y = \frac{P(x)}{Q(x)}$:
		\begin{itemize}
			\item \textbf{Bậc Tử < Bậc Mẫu:} Đồ thị luôn có đúng một TCN là trục hoành $y = 0$.
			\item \textbf{Bậc Tử = Bậc Mẫu:} TCN là đường thẳng $y = \frac{a_0}{b_0}$ (tỷ số của hai hệ số đứng trước lũy thừa bậc cao nhất ở tử và mẫu).
			\item \textbf{Bậc Tử > Bậc Mẫu:} Đồ thị không có tiệm cận ngang.
		\end{itemize}
	\end{phuongphap}
	
	\begin{vidu}[1]
		Tìm các đường tiệm cận ngang của đồ thị các hàm số sau:
		\begin{enumerate}[label=\textbf{\alph*)}]
			\item $y = f(x) = \frac{2x - 1}{x + 1}$
			\item $y = g(x) = \frac{\sqrt{4x^2 + 1}}{x - 1}$
		\end{enumerate}
	\end{vidu}
	
	\begin{loigiai}
		Theo định nghĩa, đường thẳng $y = y_0$ là tiệm cận ngang của đồ thị hàm số nếu ít nhất một trong các điều kiện sau được thỏa mãn:
		\[ \lim_{x \to +\infty} f(x) = y_0 \quad \text{hoặc} \quad \lim_{x \to -\infty} f(x) = y_0 \]
		
		\textbf{a) Khảo sát hàm số $y = \frac{2x - 1}{x + 1}$}
		
		Tập xác định của hàm số là $D = \mathbb{R} \setminus \{-1\}$. Ta tiến hành tính giới hạn của hàm số khi $x$ tiến ra vô cực:
		\begin{itemize}
			\item Khi $x \to +\infty$:
			\[ \lim_{x \to +\infty} \frac{2x - 1}{x + 1} = \lim_{x \to +\infty} \frac{x\left(2 - \frac{1}{x}\right)}{x\left(1 + \frac{1}{x}\right)} = \lim_{x \to +\infty} \frac{2 - \frac{1}{x}}{1 + \frac{1}{x}} = \frac{2 - 0}{1 + 0} = 2 \]
			\item Khi $x \to -\infty$:
			\[ \lim_{x \to -\infty} \frac{2x - 1}{x + 1} = \lim_{x \to -\infty} \frac{x\left(2 - \frac{1}{x}\right)}{x\left(1 + \frac{1}{x}\right)} = \lim_{x \to -\infty} \frac{2 - \frac{1}{x}}{1 + \frac{1}{x}} = \frac{2 - 0}{1 + 0} = 2 \]
		\end{itemize}
		Do $\lim\limits_{x \to +\infty} f(x) = \lim\limits_{x \to -\infty} f(x) = 2$, ta kết luận: Đường thẳng $y = 2$ là đường tiệm cận ngang duy nhất của đồ thị hàm số $y = \frac{2x - 1}{x + 1}$.
		
		\begin{figure}[H]
			\centering
			\begin{tikzpicture}
				\begin{axis}[
					axis lines = middle,
					xlabel = {$x$},
					ylabel = {$y$},
					xmin = -7, xmax = 5,
					ymin = -3, ymax = 7,
					xtick = {-6, -5, -4, -3, -2, -1, 1, 2, 3, 4},
					ytick = {-2, -1, 1, 2, 3, 4, 5, 6},
					grid = both,
					grid style = {dashed, gray!30},
					width = 11cm, height = 7.5cm,
					legend pos = north west
					]
					% Đồ thị nhánh bên trái (x < -1)
					\addplot [
					domain=-7:-1.2, 
					samples=100, 
					color=blue,
					thick
					] {(2*x - 1)/(x + 1)};
					
					% Đồ thị nhánh bên phải (x > -1)
					\addplot [
					domain=-0.8:5, 
					samples=100, 
					color=blue,
					thick
					] {(2*x - 1)/(x + 1)};
					
					% Tiệm cận ngang y = 2
					\addplot [
					domain=-7:5, 
					dashed, 
					color=red, 
					thick
					] {2};
					
					% Tiệm cận đứng x = -1
					\draw [dashed, color=black!50, thick] (axis cs:-1, -3) -- (axis cs:-1, 7);
					
					\node[red, above] at (axis cs:3, 2) {$y = 2$ (TCN)};
					\node[black!50, right] at (axis cs:-1, 5) {$x = -1$};
				\end{axis}
			\end{tikzpicture}
			\caption{Đồ thị hàm số $y = \frac{2x - 1}{x + 1}$ và đường tiệm cận ngang $y = 2$}
		\end{figure}
		
		\textbf{b) Khảo sát hàm số $y = \frac{\sqrt{4x^2 + 1}}{x - 1}$}
		
		Tập xác định của hàm số là $D = \mathbb{R} \setminus \{1\}$. Khi $x$ tiến ra vô cực, ta lưu ý hằng đẳng thức $\sqrt{x^2} = |x|$:
		\begin{itemize}
			\item Khi $x \to +\infty$ (do $x > 0$ nên $|x| = x$):
			\[ \lim_{x \to +\infty} \frac{\sqrt{4x^2 + 1}}{x - 1} = \lim_{x \to +\infty} \frac{\sqrt{x^2\left(4 + \frac{1}{x^2}\right)}}{x\left(1 - \frac{1}{x}\right)} = \lim_{x \to +\infty} \frac{x\sqrt{4 + \frac{1}{x^2}}}{x\left(1 - \frac{1}{x}\right)} = \lim_{x \to +\infty} \frac{\sqrt{4 + \frac{1}{x^2}}}{1 - \frac{1}{x}} = 2 \]
			\item Khi $x \to -\infty$ (do $x < 0$ nên $|x| = -x$):
			\[ \lim_{x \to -\infty} \frac{\sqrt{4x^2 + 1}}{x - 1} = \lim_{x \to -\infty} \frac{\sqrt{x^2\left(4 + \frac{1}{x^2}\right)}}{x\left(1 - \frac{1}{x}\right)} = \lim_{x \to -\infty} \frac{-x\sqrt{4 + \frac{1}{x^2}}}{x\left(1 - \frac{1}{x}\right)} = \lim_{x \to -\infty} \frac{-\sqrt{4 + \frac{1}{x^2}}}{1 - \frac{1}{x}} = -2 \]
		\end{itemize}
		Vì $\lim\limits_{x \to +\infty} g(x) = 2$ và $\lim\limits_{x \to -\infty} g(x) = -2$, đồ thị hàm số có hai đường tiệm cận ngang khác nhau là các đường thẳng:
		\[ y = 2 \quad (\text{khi } x \to +\infty) \quad \text{và} \quad y = -2 \quad (\text{khi } x \to -\infty) \]
		
		\textbf{Nhận xét quan trọng:} Đối với các hàm số chứa căn thức bậc hai dạng $\sqrt{x^2}$, học sinh cần đặc biệt lưu ý việc đổi dấu của biến $x$ khi tính giới hạn về $-\infty$. Đây là sai lầm phổ biến nhất dẫn đến việc bỏ sót đường tiệm cận ngang $y = -2$.
	\end{loigiai}
	
	
	\subsection{Đường tiệm cận đứng (TCĐ)}
	
	Đường tiệm cận đứng của đồ thị hàm số biểu diễn xu hướng đồ thị đi ra vô cực ($\pm\infty$) khi biến số $x$ tiến dần đến một giá trị hữu hạn $x_0$.
	
	\begin{phuongphap}[Định nghĩa toán học]
		Đường thẳng $x = x_0$ được gọi là đường \textbf{tiệm cận đứng} của đồ thị hàm số $y = f(x)$ nếu thỏa mãn ít nhất một trong các điều kiện sau:
		\[ 
		\lim_{x \to x_0^+} f(x) = +\infty, \quad \lim_{x \to x_0^+} f(x) = -\infty, \quad \lim_{x \to x_0^-} f(x) = +\infty, \quad \lim_{x \to x_0^-} f(x) = -\infty 
		\]
	\end{phuongphap}
	
	% === HÌNH VẼ MINH HỌA LÝ THUYẾT TỔNG QUÁT ===
	\begin{figure}[H]
		\centering
		\begin{tikzpicture}
			\begin{axis}[
				axis lines = middle,
				xlabel = {$x$},
				ylabel = {$y$},
				xmin = -1, xmax = 5,
				ymin = -6, ymax = 6,
				xtick = {2},
				xticklabels = {$x_0$},
				ytick = \empty,
				axis line style={->},
				width = 10cm, height = 7cm,
				clip = false
				]
				% Đường tiệm cận đứng x = x0 (ở đây chọn x0 = 2)
				\draw [dashed, color=black!60, thick] (axis cs:2, -6) -- (axis cs:2, 6);
				
				% Nhánh đồ thị phía bên phải (x -> 2+, y -> +inf)
				\addplot [
				domain=2.2:4.5, 
				samples=100, 
				color=blue,
				thick
				] {1.2/(x-2) + 0.5};
				
				% Nhánh đồ thị phía bên trái (x -> 2-, y -> -inf)
				\addplot [
				domain=0.2:1.8, 
				samples=100, 
				color=blue,
				thick
				] {1.2/(x-2) + 0.5};
				
				\node[blue, right] at (axis cs:3.2, 3.5) {$y = f(x)$};
				\node[black!80, right] at (axis cs:2, 5) {$x = x_0$ (Tiệm cận đứng)};
			\end{axis}
		\end{tikzpicture}
		\caption{Minh họa hình học của đường tiệm cận đứng $x = x_0$}
	\end{figure}
	
	\begin{phuongphap}[Quy tắc xác định nhanh cho hàm phân thức $y = \frac{P(x)}{Q(x)}$]
		Để tìm tiệm cận đứng của đồ thị hàm số phân thức, ta thực hiện các bước sau:
		\begin{itemize}
			\item \textbf{Bước 1:} Tìm các nghiệm của mẫu số bằng cách giải phương trình $Q(x) = 0$. Giả sử hệ phương trình có các nghiệm $x = x_1, x = x_2, \dots$
			\item \textbf{Bước 2:} Thay lần lượt các nghiệm này vào tử số $P(x)$ để kiểm chứng:
			\begin{itemize}
				\item \textbf{Trường hợp 1:} Nếu $P(x_i) \ne 0$, ta kết luận đường thẳng $x = x_i$ \textbf{chắc chắn} là đường tiệm cận đứng của đồ thị hàm số.
				\item \textbf{Trường hợp 2:} Nếu $P(x_i) = 0$ (nghiệm trùng giữa tử và mẫu), ta phải phân tích cả tử và mẫu thành nhân tử để rút gọn đại lượng trùng $(x - x_i)$. Sau đó, tính giới hạn giới hạn tại điểm gián đoạn:
				\[ \lim_{x \to x_i} \frac{P(x)}{Q(x)} \]
				Nếu giới hạn này cho ra kết quả \textbf{vô cực} ($\pm\infty$) thì đường thẳng $x = x_i$ vẫn là TCĐ. Nếu giới hạn cho ra một \textbf{số thực hữu hạn} $L$ thì đường thẳng $x = x_i$ không phải là TCĐ (đồ thị bị khuyết một điểm tại tọa độ $(x_i; L)$, gọi là điểm gián đoạn bỏ được).
			\end{itemize}
		\end{itemize}
	\end{phuongphap}
	
	\begin{vidu}[2]
		Tìm phương trình các đường tiệm cận đứng của đồ thị hàm số:
		\[ y = f(x) = \frac{x^2 - 3x + 2}{x^2 - 4} \]
	\end{vidu}
	
	\begin{loigiai}
		Tập xác định của hàm số là $D = \mathbb{R} \setminus \{-2; 2\}$.
		
		Phương trình mẫu số: $x^2 - 4 = 0 \Leftrightarrow x = 2$ hoặc $x = -2$. Ta xét từng trường hợp nghiệm:
		
		\begin{itemize}
			\item \textbf{Xét tại điểm $x = -2$:}
			Thay $x = -2$ vào tử số $P(x) = x^2 - 3x + 2$, ta được:
			\[ P(-2) = (-2)^2 - 3(-2) + 2 = 12 \ne 0 \]
			Vì tử số khác $0$ và mẫu số bằng $0$ tại $x = -2$, ta có:
			\[ \lim_{x \to -2^+} \frac{x^2 - 3x + 2}{x^2 - 4} = -\infty \quad \text{và} \quad \lim_{x \to -2^-} \frac{x^2 - 3x + 2}{x^2 - 4} = +\infty \]
			Do đó, đường thẳng $x = -2$ là đường tiệm cận đứng của đồ thị hàm số.
			
			\item \textbf{Xét tại điểm $x = 2$:}
			Thay $x = 2$ vào tử số, ta được $P(2) = 2^2 - 3(2) + 2 = 0$. Đây là trường hợp trùng nghiệm. Ta tiến hành phân tích đa thức thành nhân tử để rút gọn hàm số với $x \ne 2$:
			\[ y = \frac{x^2 - 3x + 2}{x^2 - 4} = \frac{(x - 1)(x - 2)}{(x - 2)(x + 2)} = \frac{x - 1}{x + 2} \]
			Tính giới hạn của hàm số khi $x$ dần tới $2$:
			\[ \lim_{x \to 2} y = \lim_{x \to 2} \frac{x - 1}{x + 2} = \frac{2 - 1}{2 + 2} = \frac{1}{4} \]
			Vì giới hạn này là một số thực hữu hạn ($L = \frac{1}{4}$) chứ không phải vô cực, đường thẳng $x = 2$ \textbf{không phải} là đường tiệm cận đứng của đồ thị hàm số. Tại điểm $M\left(2; \frac{1}{4}\right)$, đồ thị bị khuyết một điểm (như hình vẽ bên dưới).
		\end{itemize}
		Vậy đồ thị hàm số đã cho chỉ có duy nhất một đường tiệm cận đứng là $x = -2$.
		
		\begin{figure}[H]
			\centering
			\begin{tikzpicture}
				\begin{axis}[
					axis lines = middle,
					xlabel = {$x$},
					ylabel = {$y$},
					xmin = -8, xmax = 6,
					ymin = -4, ymax = 5,
					xtick = {-2, 2},
					ytick = {1},
					grid = both,
					grid style = {dashed, gray!30},
					width = 11cm, height = 7.5cm,
					clip = false
					]
					% Đồ thị nhánh bên trái đường tiệm cận đứng (x < -2)
					\addplot [
					domain=-8:-2.3, 
					samples=100, 
					color=blue,
					thick
					] {(x - 1)/(x + 2)};
					
					% Đồ thị nhánh bên phải đường tiệm cận đứng (x > -2) và né điểm x = 2
					\addplot [
					domain=-1.7:1.9, 
					samples=100, 
					color=blue,
					thick
					] {(x - 1)/(x + 2)};
					
					\addplot [
					domain=2.1:6, 
					samples=100, 
					color=blue,
					thick
					] {(x - 1)/(x + 2)};
					
					% Đường tiệm cận đứng x = -2
					\draw [dashed, color=red, thick] (axis cs:-2, -4) -- (axis cs:-2, 5);
					
					% Đường tiệm cận ngang y = 1 (phục vụ góc nhìn trực quan)
					\addplot [
					domain=-8:6, 
					dashed, 
					color=black!40, 
					thick
					] {1};
					
					% Điểm khuyết (hole) tại tọa độ (2, 0.25)
					\draw [blue, thick, fill=white] (axis cs:2, 0.25) circle (2.5pt);
					
					\node[red, right] at (axis cs:-2, 4) {$x = -2$ (TCĐ)};
					\node[black!60, above] at (axis cs:4, 1) {$y = 1$};
					\node[blue, above] at (axis cs:2.7, 0.3) {$\left(2; \frac{1}{4}\right)$};
				\end{axis}
			\end{tikzpicture}
			\caption{Đồ thị hàm số $y = \frac{x^2 - 3x + 2}{x^2 - 4}$ có duy nhất một TCĐ $x = -2$}
		\end{figure}
	\end{loigiai}
	
	\subsection{Đường tiệm cận xiên (TCX)}
	
	Đường tiệm cận xiên xuất hiện ở các đồ thị hàm số không có tiệm cận ngang khi $x$ tiến ra vô cực (bậc của tử số cao hơn bậc của mẫu số đúng $1$ bậc đối với hàm phân thức). Về mặt hình học, đồ thị hàm số sẽ "ôm sát" một đường thẳng nghiêng thay vì đường thẳng nằm ngang.
	
	\begin{phuongphap}[Định nghĩa toán học]
		Đường thẳng $y = ax + b$ ($a \ne 0$) được gọi là đường \textbf{tiệm cận xiên} (khi $x \to +\infty$ hoặc khi $x \to -\infty$) của đồ thị hàm số $y = f(x)$ nếu:
		\[ \lim_{x \to +\infty} [f(x) - (ax + b)] = 0 \quad \text{hoặc} \quad \lim_{x \to -\infty} [f(x) - (ax + b)] = 0 \]
	\end{phuongphap}
	
	\textbf{Bản chất hình học:} 
	Xét điểm $M(x; f(x))$ thuộc đồ thị $(C)$ và điểm $N(x; ax + b)$ thuộc đường thẳng $\Delta: y = ax + b$ có cùng hoành độ $x$. 
	\begin{itemize}
		\item Khoảng cách thẳng đứng giữa đồ thị và đường thẳng $\Delta$ chính là độ dài đoạn thẳng $MN = |f(x) - (ax + b)|$.
		\item Khi $x \to +\infty$ (hoặc $x \to -\infty$), nếu khoảng cách này tiến dần về $0$ (tức là $MN \to 0$) thì đồ thị $(C)$ và đường thẳng $\Delta$ ngày càng sát lại gần nhau. Khi đó, $\Delta$ đóng vai trò là tiệm cận xiên.
	\end{itemize}
	
	% === HÌNH VẼ MINH HỌA LÝ THUYẾT TỔNG QUÁT ===
	\begin{figure}[H]
		\centering
		\begin{tikzpicture}
			\begin{axis}[
				axis lines = middle,
				xlabel = {$x$},
				ylabel = {$y$},
				xmin = -1, xmax = 8,
				ymin = -1, ymax = 6,
				xtick = \empty,
				ytick = \empty,
				axis line style={->},
				width = 10cm, height = 6.5cm,
				clip = false
				]
				% Đường tiệm cận xiên y = ax + b (chọn y = 0.5x + 1)
				\addplot [
				domain=-1:8, 
				dashed, 
				color=black!60, 
				thick
				] {0.5*x + 1};
				
				% Đồ thị hàm số y = f(x) tiếp cận y = ax+b từ phía trên
				\addplot [
				domain=0.2:8, 
				samples=100, 
				color=blue,
				thick
				] {0.5*x + 1 + 3*exp(-0.6*x)};
				
				% Nhãn cho đồ thị và đường tiệm cận
				\node[blue, right] at (axis cs:4, 3.8) {$y = f(x)$};
				\node[black!80, below right] at (axis cs:5.5, 3.75) {$y = ax + b$};
				
				% Minh họa khoảng cách thẳng đứng MN tại điểm x_0
				\draw[dashed, black!40] (axis cs:1.5, 0) -- (axis cs:1.5, 3);
				\draw[<->, >=stealth, red, thick] (axis cs:1.5, 1.75) -- (axis cs:1.5, 2.97) node[midway, right, scale=0.8] {$MN \to 0$};
				
				\filldraw[black] (axis cs:1.5, 1.75) circle (1.5pt) node[below right, scale=0.8] {$N$};
				\filldraw[black] (axis cs:1.5, 2.97) circle (1.5pt) node[above right, scale=0.8] {$M$};
				\filldraw[black] (axis cs:1.5, 0) circle (1.2pt) node[below, scale=0.8] {$x_0$};
				
				% Mũi tên chỉ hướng biến x tiến ra dương vô cực
				\draw[->, thick, black!40] (axis cs:3, 0.4) -- (axis cs:7, 0.4) node[midway, below] {$x \to +\infty$};
			\end{axis}
		\end{tikzpicture}
		\caption{Ý nghĩa hình học của đường tiệm cận xiên $y = ax + b$ khi $x \to +\infty$}
	\end{figure}
	
	\begin{vidu}[3]
		Tìm các đường tiệm cận xiên của đồ thị các hàm số sau:
		\begin{enumerate}[label=\textbf{\alph*)}]
			\item $y = f(x) = \frac{2x^2 + x - 3}{x + 1}$
			\item $y = g(x) = \sqrt{x^2 + 2x}$
		\end{enumerate}
	\end{vidu}
	
	\begin{loigiai}
		Theo định nghĩa, đường thẳng $y = ax + b$ ($a \ne 0$) là đường tiệm cận xiên của đồ thị hàm số $y = f(x)$ khi và chỉ khi:
		\[ \lim_{x \to +\infty} [f(x) - (ax + b)] = 0 \quad \text{hoặc} \quad \lim_{x \to -\infty} [f(x) - (ax + b)] = 0 \]
		Trong đó, các hệ số $a, b$ được xác định bằng công thức giới hạn:
		\[ a = \lim_{x \to \pm\infty} \frac{f(x)}{x} \quad \text{và} \quad b = \lim_{x \to \pm\infty} [f(x) - ax] \]
		
		\textbf{a) Khảo sát hàm số $y = \frac{2x^2 + x - 3}{x + 1}$ bằng phương pháp chia đa thức}
		
		Tập xác định của hàm số là $D = \mathbb{R} \setminus \{-1\}$. Thực hiện phép chia tử số cho mẫu số, ta có:
		\[ 2x^2 + x - 3 = 2x(x + 1) - (x + 1) - 2 = (2x - 1)(x + 1) - 2 \]
		Từ đó, ta biến đổi hàm số về dạng phân tích:
		\[ y = \frac{(2x - 1)(x + 1) - 2}{x + 1} = 2x - 1 - \frac{2}{x + 1} \]
		Ta tính giới hạn của phần dư khi $x \to \pm\infty$:
		\[ \lim_{x \to \pm\infty} [y - (2x - 1)] = \lim_{x \to \pm\infty} \left( -\frac{2}{x + 1} \right) = 0 \]
		Vì giới hạn của phần dư bằng $0$ khi $x$ tiến ra vô cực, ta kết luận: Đường thẳng $y = 2x - 1$ là đường tiệm cận xiên duy nhất của đồ thị hàm số $y = \frac{2x^2 + x - 3}{x + 1}$.
		
		\begin{figure}[H]
			\centering
			\begin{tikzpicture}
				\begin{axis}[
					axis lines = middle,
					xlabel = {$x$},
					ylabel = {$y$},
					xmin = -6, xmax = 4,
					ymin = -15, ymax = 10,
					xtick = {-5, -4, -3, -2, -1, 1, 2, 3},
					ytick = {-12, -9, -6, -3, 3, 6, 9},
					grid = both,
					grid style = {dashed, gray!30},
					width = 11cm, height = 7.5cm,
					legend pos = north west
					]
					% Đồ thị nhánh trái đường tiệm cận đứng (x < -1)
					\addplot [
					domain=-6:-1.25, 
					samples=100, 
					color=blue,
					thick
					] {(2*x^2 + x - 3)/(x + 1)};
					
					% Đồ thị nhánh phải đường tiệm cận đứng (x > -1)
					\addplot [
					domain=-0.75:4, 
					samples=100, 
					color=blue,
					thick
					] {(2*x^2 + x - 3)/(x + 1)};
					
					% Tiệm cận xiên y = 2x - 1
					\addplot [
					domain=-6:4, 
					dashed, 
					color=red, 
					thick
					] {2*x - 1};
					
					% Tiệm cận đứng x = -1
					\draw [dashed, color=black!50, thick] (axis cs:-1, -15) -- (axis cs:-1, 10);
					
					\node[red, sloped, above] at (axis cs:1.8, 3.5) {$y = 2x - 1$ (TCX)};
					\node[black!60, right] at (axis cs:-1, 7) {$x = -1$};
				\end{axis}
			\end{tikzpicture}
			\caption{Đồ thị hàm số $y = \frac{2x^2 + x - 3}{x + 1}$ có TCĐ $x = -1$ và TCX $y = 2x - 1$}
		\end{figure}
		
		\textbf{b) Khảo sát hàm số $y = \sqrt{x^2 + 2x}$ bằng phương pháp giới hạn (tương đương bấm máy Casio)}
		
		Tập xác định của hàm số là $D = (-\infty; -2] \cup [0; +\infty)$. Do hàm chứa căn thức, ta tiến hành khảo sát giới hạn riêng biệt ở hai hướng $+\infty$ và $-\infty$.
		
		\begin{itemize}
			\item \textbf{Hướng thứ nhất: Khi $x \to +\infty$} (Bấm máy Casio: nhập biểu thức và gán $x = 10^5$):
			\[ a_1 = \lim_{x \to +\infty} \frac{\sqrt{x^2 + 2x}}{x} = \lim_{x \to +\infty} \frac{x\sqrt{1 + \frac{2}{x}}}{x} = \lim_{x \to +\infty} \sqrt{1 + \frac{2}{x}} = 1 \]
			\[ b_1 = \lim_{x \to +\infty} \left( \sqrt{x^2 + 2x} - x \right) = \lim_{x \to +\infty} \frac{(x^2 + 2x) - x^2}{\sqrt{x^2 + 2x} + x} = \lim_{x \to +\infty} \frac{2x}{x\left( \sqrt{1 + \frac{2}{x}} + 1 \right)} = 1 \]
			Do đó, ta có đường tiệm cận xiên thứ nhất khi $x \to +\infty$ là đường thẳng: $y = x + 1$.
			
			\item \textbf{Hướng thứ hai: Khi $x \to -\infty$} (Bấm máy Casio: nhập biểu thức và gán $x = -10^5$):
			\[ a_2 = \lim_{x \to -\infty} \frac{\sqrt{x^2 + 2x}}{x} = \lim_{x \to -\infty} \frac{-x\sqrt{1 + \frac{2}{x}}}{x} = \lim_{x \to -\infty} -\sqrt{1 + \frac{2}{x}} = -1 \]
			\[ b_2 = \lim_{x \to -\infty} \left( \sqrt{x^2 + 2x} - (-1)x \right) = \lim_{x \to -\infty} \left( \sqrt{x^2 + 2x} + x \right) \]
			Nhân liên hợp cho biểu thức trên:
			\[ b_2 = \lim_{x \to -\infty} \frac{(x^2 + 2x) - x^2}{\sqrt{x^2 + 2x} - x} = \lim_{x \to -\infty} \frac{2x}{-x\sqrt{1 + \frac{2}{x}} - x} = \lim_{x \to -\infty} \frac{2}{-\sqrt{1 + \frac{2}{x}} - 1} = -1 \]
			Do đó, ta có đường tiệm cận xiên thứ hai khi $x \to -\infty$ là đường thẳng: $y = -x - 1$.
		\end{itemize}
		
		\textbf{Bản chất toán học giải thích nhanh:} Ta có hằng đẳng thức đáng nhớ:
		\[ y = \sqrt{x^2 + 2x} = \sqrt{(x + 1)^2 - 1} \]
		Khi $x \to \pm\infty$, giá trị hằng số $1$ bên trong căn trở nên không đáng kể, do đó đồ thị hàm số sẽ tiệm cận rất sát vào đồ thị hàm trị tuyệt đối $y = |x + 1|$. 
		\begin{itemize}
			\item Khi $x \to +\infty \Rightarrow y \approx x + 1$.
			\item Khi $x \to -\infty \Rightarrow y \approx -(x + 1) = -x - 1$.
		\end{itemize}
		Đây là cách nhận diện tư duy cực kỳ sâu sắc giúp học sinh nhẩm nhanh phương trình đường tiệm cận xiên của các hàm chứa căn thức bậc hai.
	\end{loigiai}
	
	\newpage
	
	% =====================
	% PHẦN 4: BÀI TẬP
	% =====================
	\section{Bài tập luyện tập}
	
	\subsection{Câu hỏi tự luận}
	
	\subsubsection{Dạng 1: Xác định tiệm cận qua công thức của hàm số}
	
	\begin{btxanh}
		\noindent \textbf{Bài 1.} Tìm phương trình tất cả các đường tiệm cận (đứng, ngang, xiên nếu có) của đồ thị các hàm số sau:
		\begin{multicols}{3}
			\begin{enumerate}[label=\textbf{\arabic*)}]
				\item $y = \frac{4x - 1}{x - 2}$
				\item $y = \frac{2 - x}{3x + 1}$
				\item $y = \frac{5x}{x + 3}$
				\item $y = \frac{3}{2x - 4}$
				\item $y = \frac{-x + 2}{4 - x}$
				\item $y = \frac{3x - 7}{2x + 5}$
				\item $y = \frac{x^2 + 2x - 3}{x}$
				\item $y = \frac{2x^2 - x + 1}{x - 1}$
				\item $y = x - 3 + \frac{2}{x + 2}$
				\item $y = -2x + 1 - \frac{1}{x - 3}$
				\item $y = \frac{-x^2 + 3x - 1}{x + 1}$
				\item $y = \frac{3x^2}{x - 2}$
				\item $y = \frac{\sqrt{x^2 + 1}}{x}$
				\item $y = \frac{x}{\sqrt{x^2 + 4}}$
				\item $y = \frac{\sqrt{x - 1}}{x - 2}$
			\end{enumerate}
		\end{multicols}
	\end{btxanh}
	
	\vspace{0.8cm}
	
	\begin{btvang}
		\noindent \textbf{Bài 2.} Tìm phương trình tất cả các đường tiệm cận (đứng, ngang, xiên nếu có) của đồ thị các hàm số sau:
		\begin{multicols}{3}
			\begin{enumerate}[label=\textbf{\arabic*)}]
				\item $y = \frac{x^2 - 4x + 3}{x^2 - 9}$
				\item $y = \frac{x^2 - 1}{x^2 + 2x - 3}$
				\item $y = \frac{x - 2}{x^2 - 4}$
				\item $y = \frac{x^2 + x - 2}{x^2 - x - 2}$
				\item $y = \frac{2x - 4}{x^2 - 3x + 2}$
				\item $y = \frac{x^3 - 1}{x^2 - 1}$
				\item $y = \frac{x^3 - 8}{x^2 - 4}$
				\item $y = \frac{2x^2 - 3x + 1}{x^2 - 1}$
				\item $y = \frac{x^2 - x}{x^2 - 2x + 1}$
				\item $y = \frac{x^2 + 4x + 4}{x^2 - 4}$
				\item $y = \frac{3x^2 - x - 2}{x^2 - x}$
				\item $y = \frac{x^3 + 1}{x^2 - x}$
				\item $y = \frac{\sqrt{4x^2 - 1}}{x - 2}$
				\item $y = \frac{x - 1}{\sqrt{x^2 - 1}}$
				\item $y = \frac{x}{\sqrt{x^2 - 3x + 2}}$
			\end{enumerate}
		\end{multicols}
	\end{btvang}
	
	\vspace{0.8cm}
	
	\begin{btdo}
		\noindent \textbf{Bài 3.} Tìm phương trình tất cả các đường tiệm cận (đứng, ngang, xiên nếu có) của đồ thị các hàm số sau:
		\begin{multicols}{3}
			\begin{enumerate}[label=\textbf{\arabic*)}]
				\item $y = \frac{x^3 - 3x + 2}{x^2 - 2x + 1}$
				\item $y = \frac{x^3 - x^2 - x + 1}{x^2 - 3x + 2}$
				\item $y = \frac{2x^3 - x^2 - 2x + 1}{x^2 - 1}$
				\item $y = \frac{x^4 - 1}{x^3 - x^2}$
				\item $y = \frac{x^3 + x^2 - 2x}{x^2 + 3x + 2}$
				\item $y = \frac{x^3 - 8}{x^3 - 2x^2 - x + 2}$
				\item $y = \frac{x^2 - 2x + 1}{x^3 - 1}$
				\item $y = \frac{2x^3 + x^2 - 3x}{x^2 - x}$
				\item $y = \frac{x^3 + 3x^2 + 2x}{x^2 + x - 2}$
				\item $y = \frac{x^3 - 4x}{x^3 - 3x^2 + 2x}$
				\item $y = \frac{x^4 - 16}{x^3 - 2x^2}$
				\item $y = \frac{x^3 + x^2 - x - 1}{x^2 - 1}$
				\item $y = \sqrt{x^2 + 3x} - x$
				\item $y = \frac{\sqrt{x^2 - x + 1} - x}{x - 1}$
				\item $y = x + \frac{1}{\sqrt{x^2 - 1}}$
			\end{enumerate}
		\end{multicols}
	\end{btdo}
	
	\subsubsection{Dạng 2: Tìm tiệm cận thông qua Bảng biến thiên và Đồ thị (Tự luận)}
	
	\begin{btxanh}
		\textbf{Bài 4.} Tìm phương trình các đường tiệm cận của đồ thị hàm số $y=f(x)$ dựa vào bảng biến thiên được cho dưới đây:
		\vspace{0.3cm}
		
		\noindent
		\begin{minipage}{0.48\textwidth}
			\textbf{a)} 
			\begin{center}
				\begin{tikzpicture}[scale=0.85, every node/.style={scale=0.9}]
					\draw (0,0) rectangle (7,-3);
					\draw (0,-0.8) -- (7,-0.8);
					\draw (0,-1.6) -- (7,-1.6);
					\draw (1.2,0) -- (1.2,-3);
					\node at (0.6, -0.4) {$x$};
					\node at (0.6, -1.2) {$y'$};
					\node at (0.6, -2.3) {$y$};
					
					\node at (1.8, -0.4) {$-\infty$};
					\node at (4.1, -0.4) {$2$};
					\node at (6.4, -0.4) {$+\infty$};
					
					\node at (2.95, -1.2) {$-$};
					\node at (5.25, -1.2) {$-$};
					
					% Kẻ sọc đôi chuẩn xác từ dưới hàng x
					\draw (4.05, -0.8) -- (4.05, -3);
					\draw (4.15, -0.8) -- (4.15, -3);
					
					\node at (1.6, -1.9) {$1$};
					\node at (3.6, -2.7) {$-\infty$};
					\draw[->, thick] (1.9, -2.0) -- (3.3, -2.6);
					
					\node at (4.6, -1.9) {$+\infty$};
					\node at (6.6, -2.7) {$1$};
					\draw[->, thick] (4.9, -2.0) -- (6.3, -2.6);
				\end{tikzpicture}
			\end{center}
		\end{minipage}
		\hfill
		\begin{minipage}{0.48\textwidth}
			\textbf{b)}
			\begin{center}
				\begin{tikzpicture}[scale=0.85, every node/.style={scale=0.9}]
					\draw (0,0) rectangle (7,-3);
					\draw (0,-0.8) -- (7,-0.8);
					\draw (0,-1.6) -- (7,-1.6);
					\draw (1.2,0) -- (1.2,-3);
					\node at (0.6, -0.4) {$x$};
					\node at (0.6, -1.2) {$y'$};
					\node at (0.6, -2.3) {$y$};
					
					\node at (1.8, -0.4) {$-\infty$};
					\node at (4.1, -0.4) {$-1$};
					\node at (6.4, -0.4) {$+\infty$};
					
					\node at (2.95, -1.2) {$+$};
					\node at (5.25, -1.2) {$+$};
					
					% Kẻ sọc đôi chuẩn xác từ dưới hàng x
					\draw (4.05, -0.8) -- (4.05, -3);
					\draw (4.15, -0.8) -- (4.15, -3);
					
					\node at (1.6, -2.7) {$-2$};
					\node at (3.6, -1.9) {$+\infty$};
					\draw[->, thick] (1.9, -2.6) -- (3.3, -2.0);
					
					\node at (4.6, -2.7) {$-\infty$};
					\node at (6.6, -1.9) {$2$};
					\draw[->, thick] (4.9, -2.6) -- (6.3, -2.0);
				\end{tikzpicture}
			\end{center}
		\end{minipage}
		
		\vspace{0.5cm}
		\noindent
		\begin{minipage}{0.48\textwidth}
			\textbf{c)} 
			\begin{center}
				\begin{tikzpicture}[scale=0.85, every node/.style={scale=0.9}]
					\draw (0,0) rectangle (7,-3);
					\draw (0,-0.8) -- (7,-0.8);
					\draw (0,-1.6) -- (7,-1.6);
					\draw (1.2,0) -- (1.2,-3);
					\node at (0.6, -0.4) {$x$};
					\node at (0.6, -1.2) {$y'$};
					\node at (0.6, -2.3) {$y$};
					
					\node at (1.8, -0.4) {$-\infty$};
					\node at (4.1, -0.4) {$0$};
					\node at (6.4, -0.4) {$+\infty$};
					
					\node at (2.95, -1.2) {$-$};
					\node at (4.1, -1.2) {$0$};
					\node at (5.25, -1.2) {$+$};
					
					\node at (1.6, -1.9) {$3$};
					\node at (4.1, -2.7) {$-1$};
					\draw[->, thick] (2.0, -2.0) -- (3.8, -2.6);
					
					\node at (6.6, -1.9) {$3$};
					\draw[->, thick] (4.4, -2.6) -- (6.2, -2.0);
				\end{tikzpicture}
			\end{center}
		\end{minipage}
		\hfill
		\begin{minipage}{0.48\textwidth}
			\textbf{d)}
			\begin{center}
				\begin{tikzpicture}[scale=0.85, every node/.style={scale=0.9}]
					\draw (0,0) rectangle (7,-3);
					\draw (0,-0.8) -- (7,-0.8);
					\draw (0,-1.6) -- (7,-1.6);
					\draw (1.2,0) -- (1.2,-3);
					\node at (0.6, -0.4) {$x$};
					\node at (0.6, -1.2) {$y'$};
					\node at (0.6, -2.3) {$y$};
					
					\node at (1.8, -0.4) {$-\infty$};
					\node at (4.1, -0.4) {$1$};
					\node at (6.4, -0.4) {$+\infty$};
					
					\node at (2.95, -1.2) {$+$};
					\node at (5.25, -1.2) {$-$};
					
					% Kẻ sọc đôi chuẩn xác từ dưới hàng x
					\draw (4.05, -0.8) -- (4.05, -3);
					\draw (4.15, -0.8) -- (4.15, -3);
					
					\node at (1.6, -2.7) {$-\infty$};
					\node at (3.6, -1.9) {$+\infty$};
					\draw[->, thick] (1.9, -2.6) -- (3.3, -2.0);
					
					\node at (4.6, -1.9) {$+\infty$};
					\node at (6.6, -2.7) {$0$};
					\draw[->, thick] (4.9, -2.0) -- (6.3, -2.6);
				\end{tikzpicture}
			\end{center}
		\end{minipage}
		
		\vspace{0.6cm}
		\textbf{Bài 5.} Tìm phương trình các đường tiệm cận của đồ thị hàm số $y=f(x)$ được biểu diễn bằng hình vẽ sau:
		\vspace{0.3cm}
		
		\noindent
		\begin{minipage}{0.48\textwidth}
			\textbf{a)} 
			\begin{center}
				\begin{tikzpicture}
					\begin{axis}[
						width=6.5cm, height=5.5cm,
						axis lines=center,
						xmin=-4, xmax=4, ymin=-4, ymax=4,
						xtick={-1, 1}, ytick={-1, 2},
						xlabel={$x$}, ylabel={$y$},
						every axis x label/.style={at={(ticklabel* cs:1)},anchor=west},
						every axis y label/.style={at={(ticklabel* cs:1)},anchor=south},
						restrict y to domain=-10:10,
						]
						\addplot [domain=-4:-1.3, samples=50, blue, thick] {(2*x-1)/(x+1)};
						\addplot [domain=-0.7:4, samples=50, blue, thick] {(2*x-1)/(x+1)};
						\draw[dashed, red, thick] (axis cs:-1,-4) -- (axis cs:-1,4);
						\draw[dashed, red, thick] (axis cs:-4,2) -- (axis cs:4,2);
						\draw[dashed, black!60] (axis cs:1,0) -- (axis cs:1,0.5) -- (axis cs:0,0.5);
						\filldraw (axis cs:1,0.5) circle (1pt);
					\end{axis}
				\end{tikzpicture}
			\end{center}
		\end{minipage}
		\hfill
		\begin{minipage}{0.48\textwidth}
			\textbf{b)}
			\begin{center}
				\begin{tikzpicture}
					\begin{axis}[
						width=6.5cm, height=5.5cm,
						axis lines=center,
						xmin=-2, xmax=6, ymin=-4, ymax=6,
						xtick={1, 2, 3}, ytick={-1, 1},
						xlabel={$x$}, ylabel={$y$},
						every axis x label/.style={at={(ticklabel* cs:1)},anchor=west},
						every axis y label/.style={at={(ticklabel* cs:1)},anchor=south},
						restrict y to domain=-10:10,
						]
						\addplot [domain=-2:1.7, samples=50, blue, thick] {(-x+3)/(x-2)};
						\addplot [domain=2.3:6, samples=50, blue, thick] {(-x+3)/(x-2)};
						\draw[dashed, red, thick] (axis cs:2,-4) -- (axis cs:2,6);
						\draw[dashed, red, thick] (axis cs:-2,-1) -- (axis cs:6,-1);
						\draw[dashed, black!60] (axis cs:3,0) -- (axis cs:3,0) -- (axis cs:0,0);
						\filldraw (axis cs:3,0) circle (1pt);
					\end{axis}
				\end{tikzpicture}
			\end{center}
		\end{minipage}
		
		\vspace{0.5cm}
		\noindent
		\begin{minipage}{0.48\textwidth}
			\textbf{c)} 
			\begin{center}
				\begin{tikzpicture}
					\begin{axis}[
						width=6.5cm, height=5.5cm,
						axis lines=center,
						xmin=-4, xmax=4, ymin=-5, ymax=5,
						xtick={1, 2}, ytick={-2, 2},
						xlabel={$x$}, ylabel={$y$},
						every axis x label/.style={at={(ticklabel* cs:1)},anchor=west},
						every axis y label/.style={at={(ticklabel* cs:1)},anchor=south},
						restrict y to domain=-15:15,
						]
						\addplot [domain=-4:0.7, samples=50, blue, thick] {(x^2-2*x+2)/(x-1)};
						\addplot [domain=1.3:4, samples=50, blue, thick] {(x^2-2*x+2)/(x-1)};
						\draw[dashed, red, thick] (axis cs:1,-5) -- (axis cs:1,5);
						\draw[dashed, red, thick] (axis cs:-4,-5) -- (axis cs:4,3);
						\draw[dashed, black!60] (axis cs:2,0) -- (axis cs:2,2) -- (axis cs:0,2);
						\filldraw (axis cs:2,2) circle (1pt);
						\filldraw (axis cs:0,-2) circle (1pt);
					\end{axis}
				\end{tikzpicture}
			\end{center}
		\end{minipage}
		\hfill
		\begin{minipage}{0.48\textwidth}
			\textbf{d)}
			\begin{center}
				\begin{tikzpicture}
					\begin{axis}[
						width=6.5cm, height=5.5cm,
						axis lines=center,
						xmin=-5, xmax=3, ymin=-4, ymax=6,
						xtick={-2, -1}, ytick={2, 3},
						xlabel={$x$}, ylabel={$y$},
						every axis x label/.style={at={(ticklabel* cs:1)},anchor=west},
						every axis y label/.style={at={(ticklabel* cs:1)},anchor=south},
						restrict y to domain=-15:15,
						]
						\addplot [domain=-5:-1.3, samples=50, blue, thick] {(-x^2-2*x-2)/(x+1)};
						\addplot [domain=-0.7:3, samples=50, blue, thick] {(-x^2-2*x-2)/(x+1)};
						\draw[dashed, red, thick] (axis cs:-1,-4) -- (axis cs:-1,6);
						\draw[dashed, red, thick] (axis cs:-5,4) -- (axis cs:3,-4);
						\draw[dashed, black!60] (axis cs:-2,0) -- (axis cs:-2,2) -- (axis cs:0,2);
						\filldraw (axis cs:-2,2) circle (1pt);
					\end{axis}
				\end{tikzpicture}
			\end{center}
		\end{minipage}
		
	\end{btxanh}
	
	\vspace{0.8cm}
	
	\begin{btvang}
		\textbf{Bài 6.} Tìm phương trình các đường tiệm cận của đồ thị hàm số $y=f(x)$ dựa vào bảng biến thiên được cho dưới đây:
		\vspace{0.3cm}
		
		\noindent
		\begin{minipage}{0.48\textwidth}
			\textbf{a)} 
			\begin{center}
				\begin{tikzpicture}[scale=0.85, every node/.style={scale=0.9}]
					\draw (0,0) rectangle (7,-3);
					\draw (0,-0.8) -- (7,-0.8);
					\draw (0,-1.6) -- (7,-1.6);
					\draw (1.2,0) -- (1.2,-3);
					\node at (0.6, -0.4) {$x$};
					\node at (0.6, -1.2) {$y'$};
					\node at (0.6, -2.3) {$y$};
					
					\node at (1.8, -0.4) {$-\infty$};
					\node at (4.1, -0.4) {$-2$};
					\node at (6.4, -0.4) {$+\infty$};
					
					\node at (2.95, -1.2) {$+$};
					\node at (5.25, -1.2) {$-$};
					
					% Kẻ sọc đôi chuẩn xác từ dưới hàng x
					\draw (4.05, -0.8) -- (4.05, -3);
					\draw (4.15, -0.8) -- (4.15, -3);
					
					\node at (1.6, -2.7) {$-1$};
					\node at (3.6, -1.9) {$+\infty$};
					\draw[->, thick] (1.9, -2.6) -- (3.3, -2.0);
					
					\node at (4.6, -1.9) {$+\infty$};
					\node at (6.6, -2.7) {$-1$};
					\draw[->, thick] (4.9, -2.0) -- (6.3, -2.6);
				\end{tikzpicture}
			\end{center}
		\end{minipage}
		\hfill
		\begin{minipage}{0.48\textwidth}
			\textbf{b)}
			\begin{center}
				\begin{tikzpicture}[scale=0.85, every node/.style={scale=0.9}]
					\draw (0,0) rectangle (7,-3);
					\draw (0,-0.8) -- (7,-0.8);
					\draw (0,-1.6) -- (7,-1.6);
					\draw (1.2,0) -- (1.2,-3);
					\node at (0.6, -0.4) {$x$};
					\node at (0.6, -1.2) {$y'$};
					\node at (0.6, -2.3) {$y$};
					
					\node at (1.8, -0.4) {$-\infty$};
					\node at (4.1, -0.4) {$1$};
					\node at (6.4, -0.4) {$+\infty$};
					
					\node at (2.95, -1.2) {$-$};
					\node at (5.25, -1.2) {$-$};
					
					% Kẻ sọc đôi chuẩn xác từ dưới hàng x
					\draw (4.05, -0.8) -- (4.05, -3);
					\draw (4.15, -0.8) -- (4.15, -3);
					
					\node at (1.6, -1.9) {$0$};
					\node at (3.6, -2.7) {$-\infty$};
					\draw[->, thick] (1.9, -2.0) -- (3.3, -2.6);
					
					\node at (4.6, -1.9) {$+\infty$};
					\node at (6.6, -2.7) {$0$};
					\draw[->, thick] (4.9, -2.0) -- (6.3, -2.6);
				\end{tikzpicture}
			\end{center}
		\end{minipage}
		
		\vspace{0.5cm}
		\noindent
		\begin{minipage}{0.48\textwidth}
			\textbf{c)} 
			\begin{center}
				\begin{tikzpicture}[scale=0.85, every node/.style={scale=0.9}]
					\draw (0,0) rectangle (7,-3);
					\draw (0,-0.8) -- (7,-0.8);
					\draw (0,-1.6) -- (7,-1.6);
					\draw (1.2,0) -- (1.2,-3);
					\node at (0.6, -0.4) {$x$};
					\node at (0.6, -1.2) {$y'$};
					\node at (0.6, -2.3) {$y$};
					
					\node at (1.8, -0.4) {$-\infty$};
					\node at (4.1, -0.4) {$0$};
					\node at (6.4, -0.4) {$+\infty$};
					
					\node at (2.95, -1.2) {$+$};
					\node at (4.1, -1.2) {$0$};
					\node at (5.25, -1.2) {$-$};
					
					\node at (1.6, -2.7) {$2$};
					\node at (4.1, -1.9) {$5$};
					\draw[->, thick] (2.0, -2.6) -- (3.8, -2.0);
					
					\node at (6.6, -2.7) {$2$};
					\draw[->, thick] (4.4, -2.0) -- (6.2, -2.6);
				\end{tikzpicture}
			\end{center}
		\end{minipage}
		\hfill
		\begin{minipage}{0.48\textwidth}
			\textbf{d)}
			\begin{center}
				\begin{tikzpicture}[scale=0.85, every node/.style={scale=0.9}]
					\draw (0,0) rectangle (8,-3);
					\draw (0,-0.8) -- (8,-0.8);
					\draw (0,-1.6) -- (8,-1.6);
					\draw (1.2,0) -- (1.2,-3);
					\node at (0.6, -0.4) {$x$};
					\node at (0.6, -1.2) {$y'$};
					\node at (0.6, -2.3) {$y$};
					
					\node at (1.8, -0.4) {$-\infty$};
					\node at (3.5, -0.4) {$-1$};
					\node at (5.5, -0.4) {$1$};
					\node at (7.3, -0.4) {$+\infty$};
					
					\node at (2.6, -1.2) {$+$};
					\node at (4.5, -1.2) {$-$};
					\node at (6.4, -1.2) {$+$};
					
					% Kẻ sọc đôi chuẩn xác từ dưới hàng x tại x=-1
					\draw (3.45, -0.8) -- (3.45, -3);
					\draw (3.55, -0.8) -- (3.55, -3);
					
					% Kẻ sọc đôi chuẩn xác từ dưới hàng x tại x=1
					\draw (5.45, -0.8) -- (5.45, -3);
					\draw (5.55, -0.8) -- (5.55, -3);
					
					\node at (1.6, -2.7) {$-1$};
					\node at (3.2, -1.9) {$+\infty$};
					\draw[->, thick] (2.0, -2.6) -- (2.9, -2.0);
					
					\node at (3.9, -1.9) {$+\infty$};
					\node at (5.1, -2.7) {$-\infty$};
					\draw[->, thick] (4.2, -2.0) -- (4.8, -2.6);
					
					\node at (5.9, -2.7) {$-\infty$};
					\node at (7.5, -1.9) {$1$};
					\draw[->, thick] (6.2, -2.6) -- (7.1, -2.0);
				\end{tikzpicture}
			\end{center}
		\end{minipage}
		
		\vspace{0.6cm}
		\textbf{Bài 7.} Tìm phương trình các đường tiệm cận của đồ thị hàm số $y=f(x)$ được biểu diễn bằng hình vẽ sau:
		\vspace{0.3cm}
		
		\noindent
		\begin{minipage}{0.48\textwidth}
			\textbf{a)} 
			\begin{center}
				\begin{tikzpicture}
					\begin{axis}[
						width=6.5cm, height=5.5cm,
						axis lines=center,
						xmin=-4, xmax=4, ymin=-3, ymax=4,
						xtick={-1, 1}, ytick={1, 2},
						xlabel={$x$}, ylabel={$y$},
						every axis x label/.style={at={(ticklabel* cs:1)},anchor=west},
						every axis y label/.style={at={(ticklabel* cs:1)},anchor=south},
						restrict y to domain=-10:10,
						]
						\addplot [domain=-4:-1.2, samples=50, blue, thick] {1/(x^2-1) + 1};
						\addplot [domain=-0.8:0.8, samples=50, blue, thick] {1/(x^2-1) + 1};
						\addplot [domain=1.2:4, samples=50, blue, thick] {1/(x^2-1) + 1};
						\draw[dashed, red, thick] (axis cs:-1,-3) -- (axis cs:-1,4);
						\draw[dashed, red, thick] (axis cs:1,-3) -- (axis cs:1,4);
						\draw[dashed, red, thick] (axis cs:-4,1) -- (axis cs:4,1);
					\end{axis}
				\end{tikzpicture}
			\end{center}
		\end{minipage}
		\hfill
		\begin{minipage}{0.48\textwidth}
			\textbf{b)}
			\begin{center}
				\begin{tikzpicture}
					\begin{axis}[
						width=6.5cm, height=5.5cm,
						axis lines=center,
						xmin=-4, xmax=4, ymin=-5, ymax=5,
						xtick={-1, 1}, ytick={-2, 2},
						xlabel={$x$}, ylabel={$y$},
						every axis x label/.style={at={(ticklabel* cs:1)},anchor=west},
						every axis y label/.style={at={(ticklabel* cs:1)},anchor=south},
						restrict y to domain=-15:15,
						]
						\addplot [domain=-4:-0.3, samples=50, blue, thick] {(x^2+1)/x};
						\addplot [domain=0.3:4, samples=50, blue, thick] {(x^2+1)/x};
						\draw[dashed, red, thick] (axis cs:0,-5) -- (axis cs:0,5);
						\draw[dashed, red, thick] (axis cs:-4,-4) -- (axis cs:4,4);
						\draw[dashed, black!60] (axis cs:1,0) -- (axis cs:1,2) -- (axis cs:0,2);
						\filldraw (axis cs:1,2) circle (1pt);
						\draw[dashed, black!60] (axis cs:-1,0) -- (axis cs:-1,-2) -- (axis cs:0,-2);
						\filldraw (axis cs:-1,-2) circle (1pt);
					\end{axis}
				\end{tikzpicture}
			\end{center}
		\end{minipage}
		
		\vspace{0.5cm}
		\noindent
		\begin{minipage}{0.48\textwidth}
			\textbf{c)} 
			\begin{center}
				\begin{tikzpicture}
					\begin{axis}[
						width=6.5cm, height=5.5cm,
						axis lines=center,
						xmin=-4, xmax=4, ymin=-3, ymax=3,
						xtick={-1, 1}, ytick={-1, 1},
						xlabel={$x$}, ylabel={$y$},
						every axis x label/.style={at={(ticklabel* cs:1)},anchor=west},
						every axis y label/.style={at={(ticklabel* cs:1)},anchor=south},
						]
						\addplot [domain=-4:4, samples=100, blue, thick] {2*x/(x^2+1)};
						\draw[dashed, red, thick] (axis cs:-4,0) -- (axis cs:4,0);
						\draw[dashed, black!60] (axis cs:1,0) -- (axis cs:1,1) -- (axis cs:0,1);
						\filldraw (axis cs:1,1) circle (1pt);
						\draw[dashed, black!60] (axis cs:-1,0) -- (axis cs:-1,-1) -- (axis cs:0,-1);
						\filldraw (axis cs:-1,-1) circle (1pt);
					\end{axis}
				\end{tikzpicture}
			\end{center}
		\end{minipage}
		\hfill
		\begin{minipage}{0.48\textwidth}
			\textbf{d)}
			\begin{center}
				\begin{tikzpicture}
					\begin{axis}[
						width=6.5cm, height=5.5cm,
						axis lines=center,
						xmin=-3, xmax=5, ymin=-5, ymax=5,
						xtick={1, 2, 3}, ytick={-2, 2, 4},
						xlabel={$x$}, ylabel={$y$},
						every axis x label/.style={at={(ticklabel* cs:1)},anchor=west},
						every axis y label/.style={at={(ticklabel* cs:1)},anchor=south},
						restrict y to domain=-15:15,
						]
						\addplot [domain=-3:1.7, samples=50, blue, thick] {(x^2-2*x+2)/(x-2)};
						\addplot [domain=2.3:5, samples=50, blue, thick] {(x^2-2*x+2)/(x-2)};
						\draw[dashed, red, thick] (axis cs:2,-5) -- (axis cs:2,5);
						\draw[dashed, red, thick] (axis cs:-3,-3) -- (axis cs:5,5);
						\draw[dashed, black!60] (axis cs:3,0) -- (axis cs:3,5);
						\filldraw (axis cs:3,5) circle (1pt);
						\draw[dashed, black!60] (axis cs:1,0) -- (axis cs:1,-1) -- (axis cs:0,-1);
						\filldraw (axis cs:1,-1) circle (1pt);
					\end{axis}
				\end{tikzpicture}
			\end{center}
		\end{minipage}
		
	\end{btvang}
	
	\vspace{0.8cm}
	
	\begin{btdo}
		\textbf{Bài 8.} Tìm phương trình các đường tiệm cận của đồ thị hàm số $y=f(x)$ dựa vào bảng biến thiên được cho dưới đây:
		\vspace{0.3cm}
		
		\noindent
		\begin{minipage}{0.48\textwidth}
			\textbf{a)} 
			\begin{center}
				\begin{tikzpicture}[scale=0.85, every node/.style={scale=0.9}]
					\draw (0,0) rectangle (7,-3);
					\draw (0,-0.8) -- (7,-0.8);
					\draw (0,-1.6) -- (7,-1.6);
					\draw (1.2,0) -- (1.2,-3);
					\node at (0.6, -0.4) {$x$};
					\node at (0.6, -1.2) {$y'$};
					\node at (0.6, -2.3) {$y$};
					
					\node at (1.8, -0.4) {$-3$};
					\node at (4.1, -0.4) {$0$};
					\node at (6.4, -0.4) {$+\infty$};
					
					\node at (2.95, -1.2) {$-$};
					\node at (5.25, -1.2) {$+$};
					
					% Kẻ sọc đôi chuẩn xác từ dưới hàng x
					\draw (4.05, -0.8) -- (4.05, -3);
					\draw (4.15, -0.8) -- (4.15, -3);
					
					\node at (2.0, -1.9) {$4$};
					\node at (3.6, -2.7) {$-\infty$};
					\draw[->, thick] (2.3, -2.0) -- (3.3, -2.6);
					
					\node at (4.6, -2.7) {$-\infty$};
					\node at (6.6, -1.9) {$-1$};
					\draw[->, thick] (4.9, -2.6) -- (6.3, -2.0);
				\end{tikzpicture}
			\end{center}
		\end{minipage}
		\hfill
		\begin{minipage}{0.48\textwidth}
			\textbf{b)}
			\begin{center}
				\begin{tikzpicture}[scale=0.85, every node/.style={scale=0.9}]
					\draw (0,0) rectangle (8,-3);
					\draw (0,-0.8) -- (8,-0.8);
					\draw (0,-1.6) -- (8,-1.6);
					\draw (1.2,0) -- (1.2,-3);
					\node at (0.6, -0.4) {$x$};
					\node at (0.6, -1.2) {$y'$};
					\node at (0.6, -2.3) {$y$};
					
					\node at (1.8, -0.4) {$-\infty$};
					\node at (3.5, -0.4) {$-2$};
					\node at (5.5, -0.4) {$2$};
					\node at (7.3, -0.4) {$+\infty$};
					
					\node at (2.6, -1.2) {$+$};
					\node at (4.5, -1.2) {$-$};
					\node at (6.4, -1.2) {$+$};
					
					% Kẻ sọc đôi chuẩn xác từ dưới hàng x tại x=-2
					\draw (3.45, -0.8) -- (3.45, -3);
					\draw (3.55, -0.8) -- (3.55, -3);
					
					\node at (5.5, -1.2) {$0$};
					
					\node at (1.6, -2.7) {$1$};
					\node at (3.2, -1.9) {$+\infty$};
					\draw[->, thick] (2.0, -2.6) -- (2.9, -2.0);
					
					\node at (3.9, -1.9) {$+\infty$};
					\node at (5.5, -2.7) {$-3$};
					\draw[->, thick] (4.2, -2.0) -- (5.2, -2.6);
					
					\node at (7.5, -1.9) {$1$};
					\draw[->, thick] (5.8, -2.6) -- (7.1, -2.0);
				\end{tikzpicture}
			\end{center}
		\end{minipage}
		
		\vspace{0.5cm}
		\noindent
		\begin{minipage}{0.48\textwidth}
			\textbf{c)} 
			\begin{center}
				\begin{tikzpicture}[scale=0.85, every node/.style={scale=0.9}]
					\draw (0,0) rectangle (7,-3);
					\draw (0,-0.8) -- (7,-0.8);
					\draw (0,-1.6) -- (7,-1.6);
					\draw (1.2,0) -- (1.2,-3);
					\node at (0.6, -0.4) {$x$};
					\node at (0.6, -1.2) {$y'$};
					\node at (0.6, -2.3) {$y$};
					
					\node at (1.8, -0.4) {$-\infty$};
					\node at (4.1, -0.4) {$3$};
					\node at (6.4, -0.4) {$+\infty$};
					
					\node at (2.95, -1.2) {$-$};
					\node at (5.25, -1.2) {$+$};
					
					% Kẻ sọc đôi chuẩn xác từ dưới hàng x
					\draw (4.05, -0.8) -- (4.05, -3);
					\draw (4.15, -0.8) -- (4.15, -3);
					
					\node at (1.6, -1.9) {$2$};
					\node at (3.6, -2.7) {$-\infty$};
					\draw[->, thick] (2.0, -2.0) -- (3.3, -2.6);
					
					\node at (4.6, -2.7) {$5$};
					\node at (6.6, -1.9) {$+\infty$};
					\draw[->, thick] (4.9, -2.6) -- (6.2, -2.0);
				\end{tikzpicture}
			\end{center}
		\end{minipage}
		\hfill
		\begin{minipage}{0.48\textwidth}
			\textbf{d)}
			\begin{center}
				\begin{tikzpicture}[scale=0.85, every node/.style={scale=0.9}]
					\draw (0,0) rectangle (7,-3);
					\draw (0,-0.8) -- (7,-0.8);
					\draw (0,-1.6) -- (7,-1.6);
					\draw (1.2,0) -- (1.2,-3);
					\node at (0.6, -0.4) {$x$};
					\node at (0.6, -1.2) {$y'$};
					\node at (0.6, -2.3) {$y$};
					
					\node at (1.8, -0.4) {$-\infty$};
					\node at (4.1, -0.4) {$0$};
					\node at (6.4, -0.4) {$+\infty$};
					
					\node at (2.95, -1.2) {$+$};
					\node at (4.1, -1.2) {$0$};
					\node at (5.25, -1.2) {$-$};
					
					\node at (1.6, -2.7) {$-1$};
					\node at (4.1, -1.9) {$3$};
					\draw[->, thick] (2.0, -2.6) -- (3.8, -2.0);
					
					\node at (6.6, -2.7) {$1$};
					\draw[->, thick] (4.4, -2.0) -- (6.2, -2.6);
				\end{tikzpicture}
			\end{center}
		\end{minipage}
		
		\vspace{0.6cm}
		\textbf{Bài 9.} Tìm phương trình các đường tiệm cận của đồ thị hàm số $y=f(x)$ được biểu diễn bằng hình vẽ sau:
		\vspace{0.3cm}
		
		\noindent
		\begin{minipage}{0.48\textwidth}
			\textbf{a)} 
			\begin{center}
				\begin{tikzpicture}
					\begin{axis}[
						width=6.5cm, height=5.5cm,
						axis lines=center,
						xmin=-4, xmax=4, ymin=-5, ymax=3,
						xtick={-1, 1}, ytick={-2},
						xlabel={$x$}, ylabel={$y$},
						every axis x label/.style={at={(ticklabel* cs:1)},anchor=west},
						every axis y label/.style={at={(ticklabel* cs:1)},anchor=south},
						restrict y to domain=-10:10,
						]
						\addplot [domain=-4:-1.3, samples=50, blue, thick] {(-2*x)/(x+1)};
						\addplot [domain=-0.7:4, samples=50, blue, thick] {(-2*x)/(x+1)};
						\draw[dashed, red, thick] (axis cs:-1,-5) -- (axis cs:-1,3);
						\draw[dashed, red, thick] (axis cs:-4,-2) -- (axis cs:4,-2);
						\filldraw (axis cs:0,0) circle (1pt);
					\end{axis}
				\end{tikzpicture}
			\end{center}
		\end{minipage}
		\hfill
		\begin{minipage}{0.48\textwidth}
			\textbf{b)}
			\begin{center}
				\begin{tikzpicture}
					\begin{axis}[
						width=6.5cm, height=5.5cm,
						axis lines=center,
						xmin=-4, xmax=3, ymin=-5, ymax=4,
						xtick={-2, -1}, ytick={-4},
						xlabel={$x$}, ylabel={$y$},
						every axis x label/.style={at={(ticklabel* cs:1)},anchor=west},
						every axis y label/.style={at={(ticklabel* cs:1)},anchor=south},
						restrict y to domain=-15:15,
						]
						\addplot [domain=-4:-1.3, samples=50, blue, thick] {(x^2)/(x+1)};
						\addplot [domain=-0.7:3, samples=50, blue, thick] {(x^2)/(x+1)};
						\draw[dashed, red, thick] (axis cs:-1,-5) -- (axis cs:-1,4);
						\draw[dashed, red, thick] (axis cs:-4,-5) -- (axis cs:3,2);
						\draw[dashed, black!60] (axis cs:-2,0) -- (axis cs:-2,-4) -- (axis cs:0,-4);
						\filldraw (axis cs:-2,-4) circle (1pt);
						\filldraw (axis cs:0,0) circle (1pt);
					\end{axis}
				\end{tikzpicture}
			\end{center}
		\end{minipage}
		
		\vspace{0.5cm}
		\noindent
		\begin{minipage}{0.48\textwidth}
			\textbf{c)} 
			\begin{center}
				\begin{tikzpicture}
					\begin{axis}[
						width=6.5cm, height=5.5cm,
						axis lines=center,
						xmin=-2, xmax=4, ymin=-2, ymax=6,
						xtick={1}, ytick={0, 1},
						xlabel={$x$}, ylabel={$y$},
						every axis x label/.style={at={(ticklabel* cs:1)},anchor=west},
						every axis y label/.style={at={(ticklabel* cs:1)},anchor=south},
						restrict y to domain=-5:15,
						]
						\addplot [domain=-2:0.7, samples=50, blue, thick] {1/((x-1)^2)};
						\addplot [domain=1.3:4, samples=50, blue, thick] {1/((x-1)^2)};
						\draw[dashed, red, thick] (axis cs:1,-2) -- (axis cs:1,6);
						\draw[dashed, red, thick] (axis cs:-2,0) -- (axis cs:4,0);
						\draw[dashed, black!60] (axis cs:0,1) -- (axis cs:2,1) -- (axis cs:2,0);
						\filldraw (axis cs:0,1) circle (1pt);
					\end{axis}
				\end{tikzpicture}
			\end{center}
		\end{minipage}
		\hfill
		\begin{minipage}{0.48\textwidth}
			\textbf{d)}
			\begin{center}
				\begin{tikzpicture}
					\begin{axis}[
						width=6.5cm, height=5.5cm,
						axis lines=center,
						xmin=-1, xmax=6, ymin=-2, ymax=9,
						xtick={1, 2, 3}, ytick={-1, 7},
						xlabel={$x$}, ylabel={$y$},
						every axis x label/.style={at={(ticklabel* cs:1)},anchor=west},
						every axis y label/.style={at={(ticklabel* cs:1)},anchor=south},
						restrict y to domain=-15:15,
						]
						\addplot [domain=-1:1.7, samples=50, blue, thick] {(x^2-x+1)/(x-2)};
						\addplot [domain=2.3:6, samples=50, blue, thick] {(x^2-x+1)/(x-2)};
						\draw[dashed, red, thick] (axis cs:2,-2) -- (axis cs:2,9);
						\draw[dashed, red, thick] (axis cs:-1,0) -- (axis cs:6,7);
						\draw[dashed, black!60] (axis cs:3,0) -- (axis cs:3,7) -- (axis cs:0,7);
						\filldraw (axis cs:3,7) circle (1pt);
						\draw[dashed, black!60] (axis cs:1,0) -- (axis cs:1,-1) -- (axis cs:0,-1);
						\filldraw (axis cs:1,-1) circle (1pt);
					\end{axis}
				\end{tikzpicture}
			\end{center}
		\end{minipage}
		
	\end{btdo}
	
	
	
	\subsection{Câu hỏi trắc nghiệm khách quan}
	
	\noindent \textbf{Câu 1.} Cho hàm số $y = \frac{ax+b}{cx+d}$ ($c \ne 0, ad - bc \ne 0$) có đồ thị như hình vẽ dưới đây.
	\begin{center}
		\begin{tikzpicture}
			\begin{axis}[
				width=7cm, height=6cm,
				axis lines=center,
				xmin=-4, xmax=4, ymin=-4, ymax=6,
				xtick={1, 2}, ytick={-1, 2},
				xlabel={$x$}, ylabel={$y$},
				every axis x label/.style={at={(ticklabel* cs:1)},anchor=west},
				every axis y label/.style={at={(ticklabel* cs:1)},anchor=south},
				restrict y to domain=-10:10,
				]
				\addplot [domain=-4:0.7, samples=50, blue, thick] {(2*x+1)/(x-1)};
				\addplot [domain=1.3:4, samples=50, blue, thick] {(2*x+1)/(x-1)};
				\draw[dashed, black!80, thick] (axis cs:1,-4) -- (axis cs:1,6);
				\draw[dashed, black!80, thick] (axis cs:-4,2) -- (axis cs:4,2);
			\end{axis}
		\end{tikzpicture}
	\end{center}
	Đường tiệm cận đứng của đồ thị hàm số đã cho có phương trình là:
	\begin{multicols}{4}
		\begin{enumerate}[label=\textbf{\Alph*.}]
			\item $x = 1$.
			\item $x = -1$.
			\item $y = 2$.
			\item $x = 2$.
		\end{enumerate}
	\end{multicols}
	
	\vspace{0.4cm}
	\noindent \textbf{Câu 2.} Cho hàm số $y = \frac{ax+b}{cx+d}$ ($c \ne 0, ad - bc \ne 0$) có đồ thị như hình vẽ dưới đây.
	\begin{center}
		\begin{tikzpicture}
			\begin{axis}[
				width=7cm, height=6cm,
				axis lines=center,
				xmin=-6, xmax=3, ymin=-5, ymax=4,
				xtick={-2}, ytick={-1, 1},
				xlabel={$x$}, ylabel={$y$},
				every axis x label/.style={at={(ticklabel* cs:1)},anchor=west},
				every axis y label/.style={at={(ticklabel* cs:1)},anchor=south},
				restrict y to domain=-10:10,
				]
				\addplot [domain=-6:-2.3, samples=50, blue, thick] {(-x+2)/(x+2)};
				\addplot [domain=-1.7:3, samples=50, blue, thick] {(-x+2)/(x+2)};
				\draw[dashed, black!80, thick] (axis cs:-2,-5) -- (axis cs:-2,4);
				\draw[dashed, black!80, thick] (axis cs:-6,-1) -- (axis cs:3,-1);
			\end{axis}
		\end{tikzpicture}
	\end{center}
	Đường tiệm cận ngang của đồ thị hàm số đã cho có phương trình là:
	\begin{multicols}{4}
		\begin{enumerate}[label=\textbf{\Alph*.}]
			\item $x = -2$.
			\item $y = 1$.
			\item $y = -1$.
			\item $x = -1$.
		\end{enumerate}
	\end{multicols}
	
	\vspace{0.4cm}
	\noindent \textbf{Câu 3.} Cho hàm số $y = f(x)$ xác định trên $\mathbb{R} \setminus \{1\}$ và có đồ thị như hình vẽ dưới đây.
	\begin{center}
		\begin{tikzpicture}
			\begin{axis}[
				width=7cm, height=6cm,
				axis lines=center,
				xmin=-4, xmax=5, ymin=-5, ymax=6,
				xtick={1, 2}, ytick={-1, 1, 2},
				xlabel={$x$}, ylabel={$y$},
				every axis x label/.style={at={(ticklabel* cs:1)},anchor=west},
				every axis y label/.style={at={(ticklabel* cs:1)},anchor=south},
				restrict y to domain=-10:10,
				]
				\addplot [domain=-4:0.6, samples=50, blue, thick] {x - 1 - 2/(x-1)};
				\addplot [domain=1.4:5, samples=50, blue, thick] {x - 1 - 2/(x-1)};
				\draw[dashed, black!80, thick] (axis cs:1,-5) -- (axis cs:1,6);
				\draw[dashed, black!80, thick] (axis cs:-4,-5) -- (axis cs:5,4);
				\node[above, sloped] at (axis cs:-2, -3) {$y = x - 1$};
			\end{axis}
		\end{tikzpicture}
	\end{center}
	Phương trình đường tiệm cận xiên của đồ thị hàm số đã cho là:
	\begin{multicols}{4}
		\begin{enumerate}[label=\textbf{\Alph*.}]
			\item $y = x + 1$.
			\item $y = x - 1$.
			\item $y = -x + 1$.
			\item $y = x$.
		\end{enumerate}
	\end{multicols}
	
	\vspace{0.4cm}
	\noindent \textbf{Câu 4.} Cho hàm số $y = \frac{ax+b}{cx+d}$ ($c \ne 0, ad - bc \ne 0$) có đồ thị như hình vẽ dưới đây.
	\begin{center}
		\begin{tikzpicture}
			\begin{axis}[
				width=7cm, height=6cm,
				axis lines=center,
				xmin=-5, xmax=4, ymin=-5, ymax=7,
				xtick={-1, 1}, ytick={3},
				xlabel={$x$}, ylabel={$y$},
				every axis x label/.style={at={(ticklabel* cs:1)},anchor=west},
				every axis y label/.style={at={(ticklabel* cs:1)},anchor=south},
				restrict y to domain=-10:10,
				]
				\addplot [domain=-5:-1.3, samples=50, blue, thick] {(3*x)/(x+1)};
				\addplot [domain=-0.7:4, samples=50, blue, thick] {(3*x)/(x+1)};
				\draw[dashed, black!80, thick] (axis cs:-1,-5) -- (axis cs:-1,7);
				\draw[dashed, black!80, thick] (axis cs:-5,3) -- (axis cs:4,3);
			\end{axis}
		\end{tikzpicture}
	\end{center}
	Đường tiệm cận ngang của đồ thị hàm số đã cho là đường thẳng:
	\begin{multicols}{4}
		\begin{enumerate}[label=\textbf{\Alph*.}]
			\item $y = -1$.
			\item $y = 3$.
			\item $x = -1$.
			\item $x = 3$.
		\end{enumerate}
	\end{multicols}
	
	\vspace{0.4cm}
	\noindent \textbf{Câu 5.} Cho hàm số $y = \frac{ax+b}{cx+d}$ ($c \ne 0, ad - bc \ne 0$) có đồ thị như hình vẽ dưới đây.
	\begin{center}
		\begin{tikzpicture}
			\begin{axis}[
				width=7cm, height=6cm,
				axis lines=center,
				xmin=-3, xmax=6, ymin=-5, ymax=4,
				xtick={2}, ytick={-1},
				xlabel={$x$}, ylabel={$y$},
				every axis x label/.style={at={(ticklabel* cs:1)},anchor=west},
				every axis y label/.style={at={(ticklabel* cs:1)},anchor=south},
				restrict y to domain=-10:10,
				]
				\addplot [domain=-3:1.7, samples=50, blue, thick] {(-2*x+1)/(2*x-4)};
				\addplot [domain=2.3:6, samples=50, blue, thick] {(-2*x+1)/(2*x-4)};
				\draw[dashed, black!80, thick] (axis cs:2,-5) -- (axis cs:2,4);
				\draw[dashed, black!80, thick] (axis cs:-3,-1) -- (axis cs:6,-1);
			\end{axis}
		\end{tikzpicture}
	\end{center}
	Đồ thị hàm số đã cho có đường tiệm cận đứng là:
	\begin{multicols}{4}
		\begin{enumerate}[label=\textbf{\Alph*.}]
			\item $x = -1$.
			\item $y = -1$.
			\item $x = 2$.
			\item $y = 2$.
		\end{enumerate}
	\end{multicols}
	
	\vspace{0.4cm}
	\noindent \textbf{Câu 6.} Cho hàm số $y = \frac{ax+b}{cx+d}$ ($c \ne 0, ad - bc \ne 0$) có đồ thị như hình vẽ dưới đây.
	\begin{center}
		\begin{tikzpicture}
			\begin{axis}[
				width=7cm, height=6cm,
				axis lines=center,
				xmin=-5, xmax=4, ymin=-5, ymax=6,
				xtick={-1}, ytick={1, -2},
				xlabel={$x$}, ylabel={$y$},
				every axis x label/.style={at={(ticklabel* cs:1)},anchor=west},
				every axis y label/.style={at={(ticklabel* cs:1)},anchor=south},
				restrict y to domain=-10:10,
				]
				\addplot [domain=-5:-1.3, samples=50, blue, thick] {(x-2)/(x+1)};
				\addplot [domain=-0.7:4, samples=50, blue, thick] {(x-2)/(x+1)};
				\draw[dashed, black!80, thick] (axis cs:-1,-5) -- (axis cs:-1,6);
				\draw[dashed, black!80, thick] (axis cs:-5,1) -- (axis cs:4,1);
			\end{axis}
		\end{tikzpicture}
	\end{center}
	Tọa độ giao điểm $I$ của hai đường tiệm cận của đồ thị hàm số là:
	\begin{multicols}{4}
		\begin{enumerate}[label=\textbf{\Alph*.}]
			\item $I(1; -1)$.
			\item $I(-1; 1)$.
			\item $I(-1; -2)$.
			\item $I(1; 1)$.
		\end{enumerate}
	\end{multicols}
	
	
	\vspace{0.4cm}
	\noindent \textbf{Câu 7.} Cho hàm số $y = \frac{ax^2+bx+c}{mx+n}$ ($a \ne 0, m \ne 0$) với $a,b,c,m,n$ là các số thực và có đồ thị $(C)$ như hình vẽ dưới đây.
	\begin{center}
		\begin{tikzpicture}
			\begin{axis}[
				width=9cm, height=8cm,
				axis lines=center,
				xmin=-4, xmax=5, ymin=-5, ymax=6,
				xtick={1, 2}, ytick={-1, 1},
				xlabel={$x$}, ylabel={$y$},
				every axis x label/.style={at={(ticklabel* cs:1)},anchor=west},
				every axis y label/.style={at={(ticklabel* cs:1)},anchor=south},
				restrict y to domain=-10:10,
				]
				\addplot [domain=-4:0.6, samples=50, blue, thick] {x - 1 - 2/(x-1)};
				\addplot [domain=1.4:5, samples=50, blue, thick] {x - 1 - 2/(x-1)};
				\draw[dashed, black!80, thick] (axis cs:1,-5) -- (axis cs:1,6);
				\addplot [domain=-4:5, dashed, black!80, thick] {x - 1};
				\node[above, sloped] at (axis cs:-2, -3) {$y = x - 1$};
			\end{axis}
		\end{tikzpicture}
	\end{center}
	Khẳng định nào sau đây đúng?
	\begin{enumerate}[label=\textbf{\Alph*.}]
		\item Đồ thị hàm số $(C)$ có tiệm cận đứng là đường thẳng $x = 0$.
		\item Đồ thị hàm số $(C)$ có tiệm cận xiên là đường thẳng $y = 1$.
		\item Đồ thị hàm số $(C)$ có tiệm cận ngang là đường thẳng $y = x - 1$.
		\item Giao điểm của tiệm cận đứng và tiệm cận xiên có tọa độ là $(1; 0)$.
	\end{enumerate}
	
	\vspace{0.4cm}
	\noindent \textbf{Câu 8.} Cho hàm số $y = \frac{ax^2+bx+c}{mx+n}$ ($a \ne 0, m \ne 0$) với $a,b,c,m,n$ là các số thực và có đồ thị $(C)$ như hình vẽ dưới đây.
	\begin{center}
		\begin{tikzpicture}
			\begin{axis}[
				width=9cm, height=8cm,
				axis lines=center,
				xmin=-5, xmax=3, ymin=-5, ymax=4,
				xtick={-2, -1}, ytick={-1, 2},
				xlabel={$x$}, ylabel={$y$},
				every axis x label/.style={at={(ticklabel* cs:1)},anchor=west},
				every axis y label/.style={at={(ticklabel* cs:1)},anchor=south},
				restrict y to domain=-10:10,
				]
				\addplot [domain=-5:-2.4, samples=50, blue, thick] {-x - 1 - 1/(x+2)};
				\addplot [domain=-1.6:3, samples=50, blue, thick] {-x - 1 - 1/(x+2)};
				\draw[dashed, black!80, thick] (axis cs:-2,-5) -- (axis cs:-2,4);
				\addplot [domain=-5:3, dashed, black!80, thick] {-x - 1};
				\node[above, sloped, yshift=5pt] at (axis cs:2.0, -2.0) {$y = -x - 1$};
			\end{axis}
		\end{tikzpicture}
	\end{center}
	Khẳng định nào sau đây đúng?
	\begin{enumerate}[label=\textbf{\Alph*.}]
		\item Đồ thị hàm số $(C)$ có đường tiệm cận ngang.
		\item Đồ thị hàm số $(C)$ có đường tiệm cận đứng là $x = -2$ và đường tiệm cận xiên là $y = -x - 1$.
		\item Giao điểm của đường tiệm cận đứng và đường tiệm cận xiên nằm trên trục hoành.
		\item Hàm số $(C)$ có hai đường tiệm cận xiên.
	\end{enumerate}
	
	\vspace{0.4cm}
	\noindent \textbf{Câu 9.} Cho hàm số $y = \frac{ax^2+bx+c}{mx+n}$ ($a \ne 0, m \ne 0$) với $a,b,c,m,n$ là các số thực và có đồ thị $(C)$ như hình vẽ dưới đây.
	\begin{center}
		\begin{tikzpicture}
			\begin{axis}[
				width=9cm, height=8cm,
				axis lines=center,
				xmin=-2, xmax=6, ymin=-4, ymax=7,
				xtick={2, 3}, ytick={1, 2},
				xlabel={$x$}, ylabel={$y$},
				every axis x label/.style={at={(ticklabel* cs:1)},anchor=west},
				every axis y label/.style={at={(ticklabel* cs:1)},anchor=south},
				restrict y to domain=-10:10,
				]
				\addplot [domain=-2:1.6, samples=50, blue, thick] {2*x - 3 + 2/(x-2)};
				\addplot [domain=2.4:6, samples=50, blue, thick] {2*x - 3 + 2/(x-2)};
				\draw[dashed, black!80, thick] (axis cs:2,-4) -- (axis cs:2,7);
				\addplot [domain=-2:6, dashed, black!80, thick] {2*x - 3};
				\node[above, sloped] at (axis cs:5, 5) {$y = 2x - 3$};
			\end{axis}
		\end{tikzpicture}
	\end{center}
	Khẳng định nào sau đây đúng?
	\begin{enumerate}[label=\textbf{\Alph*.}]
		\item Giao điểm của tiệm cận đứng và tiệm cận xiên có tọa độ là $(2; 1)$.
		\item Đồ thị hàm số $(C)$ có tiệm cận đứng là đường thẳng $x = 2$ và tiệm cận ngang là $y = 2x - 3$.
		\item Giao điểm của hai đường tiệm cận của đồ thị hàm số nằm ở góc phần tư thứ tư của hệ tọa độ.
		\item Đồ thị hàm số $(C)$ không có đường tiệm cận đứng.
	\end{enumerate}
	
	\vspace{0.4cm}
	\noindent \textbf{Câu 10.} Cho hàm số $y = f(x)$ xác định trên $\mathbb{R} \setminus \{0\}$ và có bảng biến thiên như sau:
	\begin{center}
		\resizebox{0.7\textwidth}{!}{
			\begin{tikzpicture}[every node/.style={scale=1.1}]
				\draw[thick] (0,0) rectangle (10,-3.5);
				\draw[thick] (0,-0.8) -- (10,-0.8);
				\draw[thick] (0,-1.6) -- (10,-1.6);
				\draw[thick] (1.5,0) -- (1.5,-3.5);
				
				\node at (0.75, -0.4) {$x$};
				\node at (0.75, -1.2) {$y'$};
				\node at (0.75, -2.55) {$y$};
				
				\node at (2.2, -0.4) {$-\infty$};
				\node at (5.75, -0.4) {$0$};
				\node at (9.3, -0.4) {$+\infty$};
				
				\node at (3.975, -1.2) {$-$};
				\node at (7.525, -1.2) {$-$};
				
				% Sọc đôi từ dưới hàng x
				\draw[double, thick] (5.75, -0.8) -- (5.75, -3.5);
				
				\node at (2.2, -1.9) {$1$};
				\node at (5.2, -3.2) {$-\infty$};
				\draw[->, thick] (2.6, -2.1) -- (4.8, -3.0);
				
				\node at (6.3, -1.9) {$+\infty$};
				\node at (9.3, -3.2) {$2$};
				\draw[->, thick] (6.8, -2.1) -- (8.8, -3.0);
			\end{tikzpicture}
		}
	\end{center}
	Tổng số đường tiệm cận đứng và tiệm cận ngang của đồ thị hàm số đã cho là:
	\begin{multicols}{4}
		\begin{enumerate}[label=\textbf{\Alph*.}]
			\item $2$.
			\item $3$.
			\item $1$.
			\item $4$.
		\end{enumerate}
	\end{multicols}
	
	\vspace{0.4cm}
	\noindent \textbf{Câu 11.} Cho hàm số $y = f(x)$ xác định và liên tục trên các khoảng $(-\infty; 1)$ và $(1; +\infty)$, có bảng biến thiên như sau:
	\begin{center}
		\resizebox{0.7\textwidth}{!}{
			\begin{tikzpicture}[every node/.style={scale=1.1}]
				\draw[thick] (0,0) rectangle (10,-3.5);
				\draw[thick] (0,-0.8) -- (10,-0.8);
				\draw[thick] (0,-1.6) -- (10,-1.6);
				\draw[thick] (1.5,0) -- (1.5,-3.5);
				
				\node at (0.75, -0.4) {$x$};
				\node at (0.75, -1.2) {$y'$};
				\node at (0.75, -2.55) {$y$};
				
				\node at (2.2, -0.4) {$-\infty$};
				\node at (5.75, -0.4) {$1$};
				\node at (9.3, -0.4) {$+\infty$};
				
				\node at (3.975, -1.2) {$+$};
				\node at (7.525, -1.2) {$-$};
				
				% Sọc đôi từ dưới hàng x
				\draw[double, thick] (5.75, -0.8) -- (5.75, -3.5);
				
				\node at (2.2, -3.2) {$-\infty$};
				\node at (5.2, -1.9) {$3$};
				\draw[->, thick] (2.6, -3.0) -- (4.8, -2.1);
				
				\node at (6.3, -1.9) {$2$};
				\node at (9.3, -3.2) {$0$};
				\draw[->, thick] (6.8, -2.1) -- (8.8, -3.0);
			\end{tikzpicture}
		}
	\end{center}
	Số đường tiệm cận đứng của đồ thị hàm số đã cho là:
	\begin{multicols}{4}
		\begin{enumerate}[label=\textbf{\Alph*.}]
			\item $2$.
			\item $1$.
			\item $0$.
			\item $3$.
		\end{enumerate}
	\end{multicols}
	
	\vspace{0.4cm}
	\noindent \textbf{Câu 12.} Cho hàm số $y = f(x)$ có bảng biến thiên như sau:
	\begin{center}
		\resizebox{0.7\textwidth}{!}{
			\begin{tikzpicture}[every node/.style={scale=1.1}]
				\draw[thick] (0,0) rectangle (11,-3.5);
				\draw[thick] (0,-0.8) -- (11,-0.8);
				\draw[thick] (0,-1.6) -- (11,-1.6);
				\draw[thick] (1.5,0) -- (1.5,-3.5);
				
				\node at (0.75, -0.4) {$x$};
				\node at (0.75, -1.2) {$y'$};
				\node at (0.75, -2.55) {$y$};
				
				\node at (2.2, -0.4) {$-\infty$};
				\node at (4.5, -0.4) {$-1$};
				\node at (7.5, -0.4) {$2$};
				\node at (10.3, -0.4) {$+\infty$};
				
				\node at (3.0, -1.2) {$+$};
				\node at (6.0, -1.2) {$-$};
				\node at (8.9, -1.2) {$+$};
				
				% Sọc đôi từ dưới hàng x
				\draw[double, thick] (4.5, -0.8) -- (4.5, -3.5);
				\draw[double, thick] (7.5, -0.8) -- (7.5, -3.5);
				
				\node at (2.2, -3.2) {$1$};
				\node at (4.0, -1.9) {$+\infty$};
				\draw[->, thick] (2.5, -3.0) -- (3.6, -2.1);
				
				\node at (5.0, -1.9) {$+\infty$};
				\node at (7.0, -3.2) {$-\infty$};
				\draw[->, thick] (5.4, -2.1) -- (6.6, -3.0);
				
				\node at (8.0, -3.2) {$-\infty$};
				\node at (10.3, -1.9) {$1$};
				\draw[->, thick] (8.4, -3.0) -- (9.8, -2.1);
			\end{tikzpicture}
		}
	\end{center}
	Khẳng định nào sau đây là \textbf{sai}?
	\begin{enumerate}[label=\textbf{\Alph*.}]
		\item Đồ thị hàm số có hai đường tiệm cận đứng là $x = -1$ và $x = 2$.
		\item Đồ thị hàm số có một đường tiệm cận ngang là $y = 1$.
		\item Đồ thị hàm số có tổng cộng ba đường tiệm cận.
		\item Đồ thị hàm số có tiệm cận ngang là $x = 1$.
	\end{enumerate}
	
	\vspace{0.4cm}
	\noindent \textbf{Câu 13.} Cho hàm số $y = f(x)$ xác định trên $\mathbb{R} \setminus \{1\}$ và có bảng biến thiên như sau:
	\begin{center}
		\resizebox{0.7\textwidth}{!}{
			\begin{tikzpicture}[every node/.style={scale=1.1}]
				\draw[thick] (0,0) rectangle (10,-3.5);
				\draw[thick] (0,-0.8) -- (10,-0.8);
				\draw[thick] (0,-1.6) -- (10,-1.6);
				\draw[thick] (1.5,0) -- (1.5,-3.5);
				
				\node at (0.75, -0.4) {$x$};
				\node at (0.75, -1.2) {$y'$};
				\node at (0.75, -2.55) {$y$};
				
				\node at (2.2, -0.4) {$-\infty$};
				\node at (5.75, -0.4) {$1$};
				\node at (9.3, -0.4) {$+\infty$};
				
				\node at (3.975, -1.2) {$+$};
				\node at (7.525, -1.2) {$-$};
				
				\draw[double, thick] (5.75, -0.8) -- (5.75, -3.5);
				
				\node at (2.2, -3.2) {$2$};
				\node at (5.2, -1.9) {$5$};
				\draw[->, thick] (2.6, -3.0) -- (4.8, -2.1);
				
				\node at (6.3, -1.9) {$+\infty$};
				\node at (9.3, -3.2) {$-\infty$};
				\draw[->, thick] (6.8, -2.1) -- (8.8, -3.0);
			\end{tikzpicture}
		}
	\end{center}
	Tổng số tiệm cận đứng và tiệm cận ngang của đồ thị hàm số đã cho là:
	\begin{multicols}{4}
		\begin{enumerate}[label=\textbf{\Alph*.}]
			\item $1$.
			\item $2$.
			\item $3$.
			\item $0$.
		\end{enumerate}
	\end{multicols}
	
	\vspace{0.4cm}
	\noindent \textbf{Câu 14.} Cho hàm số đa thức bậc ba $y = f(x)$ có bảng biến thiên như sau:
	\begin{center}
		\resizebox{0.7\textwidth}{!}{
			\begin{tikzpicture}[every node/.style={scale=1.1}]
				\draw[thick] (0,0) rectangle (10,-3.5);
				\draw[thick] (0,-0.8) -- (10,-0.8);
				\draw[thick] (0,-1.6) -- (10,-1.6);
				\draw[thick] (1.5,0) -- (1.5,-3.5);
				
				\node at (0.75, -0.4) {$x$};
				\node at (0.75, -1.2) {$y'$};
				\node at (0.75, -2.55) {$y$};
				
				\node at (2.2, -0.4) {$-\infty$};
				\node at (4.5, -0.4) {$-1$};
				\node at (7.0, -0.4) {$1$};
				\node at (9.3, -0.4) {$+\infty$};
				
				\node at (3.35, -1.2) {$+$};
				\node at (4.5, -1.2) {$0$};
				\node at (5.75, -1.2) {$-$};
				\node at (7.0, -1.2) {$0$};
				\node at (8.15, -1.2) {$+$};
				
				\draw (4.5, -0.8) -- (4.5, -1.6);
				\draw (7.0, -0.8) -- (7.0, -1.6);
				
				\node at (2.2, -3.2) {$-\infty$};
				\node at (4.5, -1.9) {$3$};
				\draw[->, thick] (2.6, -3.0) -- (4.1, -2.1);
				
				\node at (7.0, -3.2) {$-1$};
				\draw[->, thick] (4.9, -2.1) -- (6.6, -3.0);
				
				\node at (9.3, -1.9) {$+\infty$};
				\draw[->, thick] (7.4, -3.0) -- (8.9, -2.1);
			\end{tikzpicture}
		}
	\end{center}
	Đồ thị hàm số đã cho có bao nhiêu đường tiệm cận?
	\begin{multicols}{4}
		\begin{enumerate}[label=\textbf{\Alph*.}]
			\item $2$.
			\item $1$.
			\item $3$.
			\item $0$.
		\end{enumerate}
	\end{multicols}
	
	\vspace{0.4cm}
	\noindent \textbf{Câu 15.} Cho hàm số $y = f(x)$ xác định trên $\mathbb{R} \setminus \{3\}$ và có bảng biến thiên như sau:
	\begin{center}
		\resizebox{0.7\textwidth}{!}{
			\begin{tikzpicture}[every node/.style={scale=1.1}]
				\draw[thick] (0,0) rectangle (10,-3.5);
				\draw[thick] (0,-0.8) -- (10,-0.8);
				\draw[thick] (0,-1.6) -- (10,-1.6);
				\draw[thick] (1.5,0) -- (1.5,-3.5);
				
				\node at (0.75, -0.4) {$x$};
				\node at (0.75, -1.2) {$y'$};
				\node at (0.75, -2.55) {$y$};
				
				\node at (2.2, -0.4) {$-\infty$};
				\node at (5.75, -0.4) {$3$};
				\node at (9.3, -0.4) {$+\infty$};
				
				\node at (3.975, -1.2) {$-$};
				\node at (7.525, -1.2) {$-$};
				
				\draw[double, thick] (5.75, -0.8) -- (5.75, -3.5);
				
				\node at (2.2, -1.9) {$-2$};
				\node at (5.2, -3.2) {$-\infty$};
				\draw[->, thick] (2.6, -2.1) -- (4.8, -3.0);
				
				\node at (6.3, -1.9) {$+\infty$};
				\node at (9.3, -3.2) {$-2$};
				\draw[->, thick] (6.8, -2.1) -- (8.8, -3.0);
			\end{tikzpicture}
		}
	\end{center}
	Phương trình các đường tiệm cận của đồ thị hàm số là:
	\begin{multicols}{2}
		\begin{enumerate}[label=\textbf{\Alph*.}]
			\item $x = -2$ và $y = 3$.
			\item $x = 3$ và $y = -2$.
			\item $x = 3$ và $y = 2$.
			\item Không có tiệm cận.
		\end{enumerate}
	\end{multicols}
	
	\vspace{0.4cm}
	\noindent \textbf{Câu 16.} Cho hàm số $y = f(x)$ xác định trên $(-\infty; -2) \cup (2; +\infty)$ và có bảng biến thiên như sau:
	\begin{center}
		\resizebox{0.7\textwidth}{!}{
			\begin{tikzpicture}[every node/.style={scale=1.1}]
				\draw[thick] (0,0) rectangle (11,-3.5);
				\draw[thick] (0,-0.8) -- (11,-0.8);
				\draw[thick] (0,-1.6) -- (11,-1.6);
				\draw[thick] (1.5,0) -- (1.5,-3.5);
				
				\node at (0.75, -0.4) {$x$};
				\node at (0.75, -1.2) {$y'$};
				\node at (0.75, -2.55) {$y$};
				
				\node at (2.2, -0.4) {$-\infty$};
				\node at (4.5, -0.4) {$-2$};
				\node at (7.5, -0.4) {$2$};
				\node at (10.3, -0.4) {$+\infty$};
				
				\node at (3.35, -1.2) {$-$};
				\node at (6.0, -1.2) {}; 
				\node at (8.9, -1.2) {$+$};
				
				\draw[double, thick] (4.5, -0.8) -- (4.5, -3.5);
				\draw[double, thick] (7.5, -0.8) -- (7.5, -3.5);
				
				% Vùng không xác định được tô màu xám nhạt (thay cho gạch chéo để tránh lỗi package)
				\fill[black!15] (4.55, -0.8) rectangle (7.45, -3.5);
				
				\node at (2.2, -1.9) {$-1$};
				\node at (4.0, -3.2) {$-\infty$};
				\draw[->, thick] (2.5, -2.1) -- (3.6, -3.0);
				
				\node at (8.0, -3.2) {$0$};
				\node at (10.3, -1.9) {$1$};
				\draw[->, thick] (8.4, -3.0) -- (9.8, -2.1);
			\end{tikzpicture}
		}
	\end{center}
	Đồ thị hàm số đã cho có tổng cộng bao nhiêu đường tiệm cận đứng và tiệm cận ngang?
	\begin{multicols}{4}
		\begin{enumerate}[label=\textbf{\Alph*.}]
			\item $2$.
			\item $3$.
			\item $4$.
			\item $1$.
		\end{enumerate}
	\end{multicols}
	
	\vspace{0.4cm}
	\noindent \textbf{Câu 17.} Cho hàm số $y = f(x)$ có bảng biến thiên như sau:
	\begin{center}
		\resizebox{0.7\textwidth}{!}{
			\begin{tikzpicture}[every node/.style={scale=1.1}]
				\draw[thick] (0,0) rectangle (10,-3.5);
				\draw[thick] (0,-0.8) -- (10,-0.8);
				\draw[thick] (0,-1.6) -- (10,-1.6);
				\draw[thick] (1.5,0) -- (1.5,-3.5);
				
				\node at (0.75, -0.4) {$x$};
				\node at (0.75, -1.2) {$y'$};
				\node at (0.75, -2.55) {$y$};
				
				\node at (2.2, -0.4) {$-\infty$};
				\node at (5.75, -0.4) {$-2$};
				\node at (9.3, -0.4) {$+\infty$};
				
				\node at (3.975, -1.2) {$-$};
				\node at (7.525, -1.2) {$+$};
				
				\draw[double, thick] (5.75, -0.8) -- (5.75, -3.5);
				
				\node at (2.2, -1.9) {$-1$};
				\node at (5.2, -3.2) {$-\infty$};
				\draw[->, thick] (2.6, -2.1) -- (4.8, -3.0);
				
				\node at (6.3, -3.2) {$-\infty$};
				\node at (9.3, -1.9) {$3$};
				\draw[->, thick] (6.8, -3.0) -- (8.8, -2.1);
			\end{tikzpicture}
		}
	\end{center}
	Đồ thị hàm số đã cho có các đường tiệm cận là:
	\begin{enumerate}[label=\textbf{\Alph*.}]
		\item $x = -2, y = -1$ và $y = 3$.
		\item $x = -2$ và $y = 3$.
		\item $x = -1, x = 3$ và $y = -2$.
		\item $y = -2, y = -1$ và $y = 3$.
	\end{enumerate}
	
	\vspace{0.4cm}
	\noindent \textbf{Câu 18.} (Trích từ đề bài gốc của giáo viên) Cho hàm số $y = f(x)$ có bảng biến thiên như sau:
	\begin{center}
		\resizebox{0.7\textwidth}{!}{
			\begin{tikzpicture}[every node/.style={scale=1.1}]
				\draw[thick] (0,0) rectangle (10,-3.5);
				\draw[thick] (0,-0.8) -- (10,-0.8);
				\draw[thick] (0,-1.6) -- (10,-1.6);
				\draw[thick] (1.5,0) -- (1.5,-3.5);
				
				\node at (0.75, -0.4) {$x$};
				\node at (0.75, -1.2) {$y'$};
				\node at (0.75, -2.55) {$y$};
				
				\node at (2.2, -0.4) {$-\infty$};
				\node at (5.0, -0.4) {$0$};
				\node at (7.5, -0.4) {$3$};
				\node at (9.3, -0.4) {$+\infty$};
				
				\node at (3.6, -1.2) {$-$};
				\node at (6.25, -1.2) {$-$};
				\node at (8.4, -1.2) {$+$};
				
				\draw[double, thick] (5.0, -0.8) -- (5.0, -3.5);
				
				\draw (7.5, -0.8) -- (7.5, -1.6);
				\node at (7.5, -1.2) {$0$};
				
				\node at (2.2, -1.9) {$1$};
				\node at (4.5, -3.2) {$-\infty$};
				\draw[->, thick] (2.6, -2.1) -- (4.1, -3.0);
				
				\node at (5.5, -1.9) {$2$};
				\node at (7.5, -3.2) {$-3$};
				\draw[->, thick] (5.8, -2.1) -- (7.1, -3.0);
				
				\node at (9.3, -1.9) {$3$};
				\draw[->, thick] (7.9, -3.0) -- (8.9, -2.1);
			\end{tikzpicture}
		}
	\end{center}
	Tổng số tiệm cận đứng và tiệm cận ngang của đồ thị hàm số đã cho là:
	\begin{multicols}{4}
		\begin{enumerate}[label=\textbf{\Alph*.}]
			\item $2$.
			\item $3$.
			\item $4$.
			\item $1$.
		\end{enumerate}
	\end{multicols}
	
	\vspace{0.4cm}
	\noindent \textbf{Câu 19.} Cho hàm số $y = f(x)$ xác định trên $\mathbb{R} \setminus \{-1; 2\}$ và có bảng biến thiên như sau:
	\begin{center}
		\resizebox{0.7\textwidth}{!}{
			\begin{tikzpicture}[every node/.style={scale=1.1}]
				\draw[thick] (0,0) rectangle (11,-3.5);
				\draw[thick] (0,-0.8) -- (11,-0.8);
				\draw[thick] (0,-1.6) -- (11,-1.6);
				\draw[thick] (1.5,0) -- (1.5,-3.5);
				
				\node at (0.75, -0.4) {$x$};
				\node at (0.75, -1.2) {$y'$};
				\node at (0.75, -2.55) {$y$};
				
				\node at (2.2, -0.4) {$-\infty$};
				\node at (4.5, -0.4) {$-1$};
				\node at (7.5, -0.4) {$2$};
				\node at (10.3, -0.4) {$+\infty$};
				
				\node at (3.35, -1.2) {$+$};
				\node at (6.0, -1.2) {$-$};
				\node at (8.9, -1.2) {$-$};
				
				\draw[double, thick] (4.5, -0.8) -- (4.5, -3.5);
				\draw[double, thick] (7.5, -0.8) -- (7.5, -3.5);
				
				\node at (2.2, -3.2) {$+\infty$};
				\node at (4.0, -1.9) {$2$};
				\draw[->, thick] (2.5, -3.0) -- (3.6, -2.1);
				
				\node at (5.0, -1.9) {$+\infty$};
				\node at (7.0, -3.2) {$-\infty$};
				\draw[->, thick] (5.4, -2.1) -- (6.6, -3.0);
				
				\node at (8.0, -1.9) {$+\infty$};
				\node at (10.3, -3.2) {$0$};
				\draw[->, thick] (8.4, -2.1) -- (9.8, -3.0);
			\end{tikzpicture}
		}
	\end{center}
	Số lượng đường tiệm cận đứng và tiệm cận ngang của đồ thị hàm số lần lượt là:
	\begin{multicols}{2}
		\begin{enumerate}[label=\textbf{\Alph*.}]
			\item $2$ tiệm cận đứng và $1$ tiệm cận ngang.
			\item $1$ tiệm cận đứng và $1$ tiệm cận ngang.
			\item $2$ tiệm cận đứng và $0$ tiệm cận ngang.
			\item $1$ tiệm cận đứng và $2$ tiệm cận ngang.
		\end{enumerate}
	\end{multicols}
	
	
	\vspace{0.4cm}
	\noindent \textbf{Câu 20.} Tiệm cận đứng của đồ thị hàm số $y = \frac{4x - 1}{x + 3}$ là đường thẳng có phương trình:
	\begin{multicols}{4}
		\begin{enumerate}[label=\textbf{\Alph*.}]
			\item $x = -3$.
			\item $x = 4$.
			\item $x = 3$.
			\item $y = -3$.
		\end{enumerate}
	\end{multicols}
	
	\vspace{0.4cm}
	\noindent \textbf{Câu 21.} Tiệm cận ngang của đồ thị hàm số $y = \frac{1 - 2x}{x + 2}$ là đường thẳng có phương trình:
	\begin{multicols}{4}
		\begin{enumerate}[label=\textbf{\Alph*.}]
			\item $y = 1$.
			\item $y = -2$.
			\item $x = -2$.
			\item $y = -\frac{1}{2}$.
		\end{enumerate}
	\end{multicols}
	
	\vspace{0.4cm}
	\noindent \textbf{Câu 22.} Đồ thị hàm số $y = \frac{5x}{2 - x}$ có đường tiệm cận đứng là:
	\begin{multicols}{4}
		\begin{enumerate}[label=\textbf{\Alph*.}]
			\item $x = 2$.
			\item $x = -2$.
			\item $y = -5$.
			\item $x = \frac{5}{2}$.
		\end{enumerate}
	\end{multicols}
	
	\vspace{0.4cm}
	\noindent \textbf{Câu 23.} Phương trình tiệm cận ngang của đồ thị hàm số $y = \frac{3x + 4}{2x - 1}$ là:
	\begin{multicols}{4}
		\begin{enumerate}[label=\textbf{\Alph*.}]
			\item $y = 3$.
			\item $y = \frac{3}{2}$.
			\item $x = \frac{1}{2}$.
			\item $y = -4$.
		\end{enumerate}
	\end{multicols}
	
	\vspace{0.4cm}
	\noindent \textbf{Câu 24.} Giao điểm của hai đường tiệm cận của đồ thị hàm số $y = \frac{2x - 5}{x + 1}$ có tọa độ là:
	\begin{multicols}{4}
		\begin{enumerate}[label=\textbf{\Alph*.}]
			\item $(-1; 2)$.
			\item $(2; -1)$.
			\item $(-1; -5)$.
			\item $(1; 2)$.
		\end{enumerate}
	\end{multicols}
	
	\vspace{0.4cm}
	\noindent \textbf{Câu 25.} Đồ thị hàm số nào dưới đây có tiệm cận đứng là đường thẳng $x = 1$?
	\begin{multicols}{4}
		\begin{enumerate}[label=\textbf{\Alph*.}]
			\item $y = \frac{2x}{x + 1}$.
			\item $y = \frac{x - 1}{x + 2}$.
			\item $y = \frac{3x + 1}{1 - x}$.
			\item $y = \frac{x}{x^2 + 1}$.
		\end{enumerate}
	\end{multicols}
	
	\vspace{0.4cm}
	\noindent \textbf{Câu 26.} Đồ thị hàm số nào dưới đây có tiệm cận ngang là đường thẳng $y = -1$?
	\begin{multicols}{4}
		\begin{enumerate}[label=\textbf{\Alph*.}]
			\item $y = \frac{-x + 2}{x - 1}$.
			\item $y = \frac{x + 2}{-x^2 + 1}$.
			\item $y = \frac{2x - 1}{x + 1}$.
			\item $y = \frac{-x^2 + 1}{x - 1}$.
		\end{enumerate}
	\end{multicols}
	
	\vspace{0.4cm}
	\noindent \textbf{Câu 27.} Phương trình đường tiệm cận xiên của đồ thị hàm số $y = \frac{x^2 - 3x + 2}{x + 1}$ là:
	\begin{multicols}{4}
		\begin{enumerate}[label=\textbf{\Alph*.}]
			\item $y = x - 4$.
			\item $y = x - 3$.
			\item $y = x - 2$.
			\item $y = x + 1$.
		\end{enumerate}
	\end{multicols}
	
	\vspace{0.4cm}
	\noindent \textbf{Câu 28.} Tiệm cận đứng và tiệm cận xiên của đồ thị hàm số $y = \frac{x^2 - x - 1}{x - 2}$ lần lượt là:
	\begin{multicols}{2}
		\begin{enumerate}[label=\textbf{\Alph*.}]
			\item $x = 2$ và $y = x - 1$.
			\item $x = -2$ và $y = x + 1$.
			\item $x = 2$ và $y = x + 1$.
			\item $x = 2$ và $y = x - 2$.
		\end{enumerate}
	\end{multicols}
	
	\vspace{0.4cm}
	\noindent \textbf{Câu 29.} Đồ thị hàm số $y = x - 2 + \frac{3}{x - 1}$ có đường tiệm cận xiên là:
	\begin{multicols}{4}
		\begin{enumerate}[label=\textbf{\Alph*.}]
			\item $y = x$.
			\item $y = x - 2$.
			\item $y = x - 1$.
			\item $y = x + 3$.
		\end{enumerate}
	\end{multicols}
	
	\vspace{0.4cm}
	\noindent \textbf{Câu 30.} Đồ thị hàm số $y = \frac{2x^2 + x - 3}{x}$ có đường tiệm cận xiên đi qua điểm nào dưới đây?
	\begin{multicols}{4}
		\begin{enumerate}[label=\textbf{\Alph*.}]
			\item $(0; 0)$.
			\item $(1; 3)$.
			\item $(2; 5)$.
			\item $(1; 2)$.
		\end{enumerate}
	\end{multicols}
	
	\vspace{0.4cm}
	\noindent \textbf{Câu 31.} Đồ thị hàm số $y = \frac{x^2 - 4}{x^2 - 1}$ có bao nhiêu đường tiệm cận đứng?
	\begin{multicols}{4}
		\begin{enumerate}[label=\textbf{\Alph*.}]
			\item $0$.
			\item $1$.
			\item $2$.
			\item $3$.
		\end{enumerate}
	\end{multicols}
	
	\vspace{0.4cm}
	\noindent \textbf{Câu 32.} Số đường tiệm cận ngang của đồ thị hàm số $y = \frac{2x - 1}{\sqrt{x^2 + 1}}$ là:
	\begin{multicols}{4}
		\begin{enumerate}[label=\textbf{\Alph*.}]
			\item $0$.
			\item $1$.
			\item $2$.
			\item $3$.
		\end{enumerate}
	\end{multicols}
	
	\vspace{0.4cm}
	\noindent \textbf{Câu 33.} Đồ thị hàm số $y = \frac{\sqrt{x^2 - x}}{x - 1}$ có tiệm cận đứng là:
	\begin{multicols}{4}
		\begin{enumerate}[label=\textbf{\Alph*.}]
			\item $x = 1$.
			\item $x = 0$.
			\item $y = 1$.
			\item Không có tiệm cận đứng.
		\end{enumerate}
	\end{multicols}
	
	\vspace{0.4cm}
	\noindent \textbf{Câu 34.} Cho hàm số $y = \frac{x - 2}{x^2 - 4}$. Khẳng định nào sau đây là \textbf{đúng}?
	\begin{enumerate}[label=\textbf{\Alph*.}]
		\item Đồ thị hàm số có hai tiệm cận đứng là $x = 2$ và $x = -2$.
		\item Đồ thị hàm số có một tiệm cận đứng là $x = -2$ và một tiệm cận ngang là $y = 0$.
		\item Đồ thị hàm số có một tiệm cận đứng là $x = 2$ và một tiệm cận ngang là $y = 0$.
		\item Đồ thị hàm số không có tiệm cận ngang.
	\end{enumerate}
	
	\vspace{0.4cm}
	\noindent \textbf{Câu 35.} Đường thẳng nào dưới đây là tiệm cận ngang của đồ thị hàm số $y = \frac{3x^2 + 1}{x^2 - 2x + 1}$?
	\begin{multicols}{4}
		\begin{enumerate}[label=\textbf{\Alph*.}]
			\item $y = 1$.
			\item $y = 3$.
			\item $x = 1$.
			\item $y = 0$.
		\end{enumerate}
	\end{multicols}
	
	\vspace{0.4cm}
	\noindent \textbf{Câu 36.} Đồ thị hàm số $y = \frac{x^3 - 1}{x^2 - 3x + 2}$ có đường tiệm cận xiên là:
	\begin{multicols}{4}
		\begin{enumerate}[label=\textbf{\Alph*.}]
			\item $y = x + 3$.
			\item $y = x - 3$.
			\item $y = x + 2$.
			\item $y = x - 2$.
		\end{enumerate}
	\end{multicols}
	
	\vspace{0.4cm}
	\noindent \textbf{Câu 37.} Tọa độ giao điểm của tiệm cận đứng và tiệm cận xiên của đồ thị hàm số $y = \frac{x^2 + 2x + 2}{x + 1}$ là:
	\begin{multicols}{4}
		\begin{enumerate}[label=\textbf{\Alph*.}]
			\item $(-1; 0)$.
			\item $(1; 2)$.
			\item $(-1; 1)$.
			\item $(0; 1)$.
		\end{enumerate}
	\end{multicols}
	
	\vspace{0.4cm}
	\noindent \textbf{Câu 38.} Hàm số nào dưới đây có đồ thị \textbf{không} có tiệm cận đứng?
	\begin{multicols}{4}
		\begin{enumerate}[label=\textbf{\Alph*.}]
			\item $y = \frac{x - 1}{x^2 - 1}$.
			\item $y = \frac{x^2 + 1}{x - 1}$.
			\item $y = \frac{x}{x^2 + 1}$.
			\item $y = \frac{x + 1}{x^2 - x}$.
		\end{enumerate}
	\end{multicols}
	
	\vspace{0.4cm}
	\noindent \textbf{Câu 39.} Tổng số đường tiệm cận (đứng và ngang) của đồ thị hàm số $y = \frac{x + 1}{\sqrt{x^2 - 4}}$ là:
	\begin{multicols}{4}
		\begin{enumerate}[label=\textbf{\Alph*.}]
			\item $2$.
			\item $3$.
			\item $4$.
			\item $1$.
		\end{enumerate}
	\end{multicols}
	
	
	\vspace{0.4cm}
	\noindent \textbf{Câu 40.} Xác định tọa độ giao điểm $I$ của đường tiệm cận đứng và đường tiệm cận xiên của đồ thị hàm số $y = \frac{3x^2 - 5x + 4}{x - 2}$.
	\begin{multicols}{4}
		\begin{enumerate}[label=\textbf{\Alph*.}]
			\item $I(2; 7)$.
			\item $I(2; 1)$.
			\item $I(2; -5)$.
			\item $I(2; 5)$.
		\end{enumerate}
	\end{multicols}
	
	\vspace{0.4cm}
	\noindent \textbf{Câu 41.} Số đường tiệm cận đứng của đồ thị hàm số $y = \frac{\sqrt{x+9} - 3}{x^2 + x}$ là:
	\begin{multicols}{4}
		\begin{enumerate}[label=\textbf{\Alph*.}]
			\item $2$.
			\item $1$.
			\item $3$.
			\item $0$.
		\end{enumerate}
	\end{multicols}
	
	\vspace{0.4cm}
	\noindent \textbf{Câu 42.} Đồ thị hàm số $y = \frac{x - 2}{\sqrt{x^2 - 4}}$ có tất cả bao nhiêu đường tiệm cận đứng và tiệm cận ngang?
	\begin{multicols}{4}
		\begin{enumerate}[label=\textbf{\Alph*.}]
			\item $4$.
			\item $2$.
			\item $3$.
			\item $1$.
		\end{enumerate}
	\end{multicols}
	
	\vspace{0.4cm}
	\noindent \textbf{Câu 43.} Tọa độ giao điểm $I$ của các đường tiệm cận của đồ thị hàm số $y = \frac{x^2 - x + 1}{x - 1}$ là:
	\begin{multicols}{4}
		\begin{enumerate}[label=\textbf{\Alph*.}]
			\item $I(1; 1)$.
			\item $I(1; 0)$.
			\item $I(1; -1)$.
			\item $I(0; 1)$.
		\end{enumerate}
	\end{multicols}
	
	\vspace{0.4cm}
	\noindent \textbf{Câu 44.} Số đường tiệm cận đứng của đồ thị hàm số $y = \frac{\sqrt{x+1} - 1}{x^2 - x}$ là:
	\begin{multicols}{4}
		\begin{enumerate}[label=\textbf{\Alph*.}]
			\item $0$.
			\item $1$.
			\item $2$.
			\item $3$.
		\end{enumerate}
	\end{multicols}
	
	\vspace{0.4cm}
	\noindent \textbf{Câu 45.} Tổng số đường tiệm cận đứng và tiệm cận ngang của đồ thị hàm số $y = \frac{2x - 1}{\sqrt{x^2 - 9}}$ là:
	\begin{multicols}{4}
		\begin{enumerate}[label=\textbf{\Alph*.}]
			\item $2$.
			\item $3$.
			\item $4$.
			\item $1$.
		\end{enumerate}
	\end{multicols}
	
	
	
	\subsection{Câu hỏi đúng,sai}
	
	\noindent \textbf{Câu 36.} Cho hàm số $y = \frac{ax^2+bx+c}{mx+n}$ ($a \ne 0, m \ne 0$) với $a,b,c,m,n$ là các số thực và có đồ thị $(C)$ như hình vẽ dưới đây.
	\begin{center}
		\begin{tikzpicture}
			\begin{axis}[
				width=9cm, height=8cm,
				axis lines=center,
				xmin=-2, xmax=6, ymin=-2, ymax=8,
				xtick={1, 2, 3}, ytick={1, 5},
				xlabel={$x$}, ylabel={$y$},
				every axis x label/.style={at={(ticklabel* cs:1)},anchor=west},
				every axis y label/.style={at={(ticklabel* cs:1)},anchor=south},
				restrict y to domain=-5:15,
				]
				\addplot [domain=-2:1.7, samples=50, blue, thick] {(x^2-x-1)/(x-2)};
				\addplot [domain=2.3:6, samples=50, blue, thick] {(x^2-x-1)/(x-2)};
				\draw[dashed, red, thick] (axis cs:2,-2) -- (axis cs:2,8);
				\draw[dashed, red, thick] (axis cs:-2,-1) -- (axis cs:6,7);
				\draw[dashed, black!60] (axis cs:1,0) -- (axis cs:1,1) -- (axis cs:0,1);
				\draw[dashed, black!60] (axis cs:3,0) -- (axis cs:3,5) -- (axis cs:0,5);
				\filldraw (axis cs:1,1) circle (1.5pt);
				\filldraw (axis cs:3,5) circle (1.5pt);
				\node[above, sloped] at (axis cs:4.5, 3.5) {$y = x + 1$};
			\end{axis}
		\end{tikzpicture}
	\end{center}
	Các mệnh đề sau đúng hay sai?
	\begin{enumerate}[label=\textbf{\alph*)}]
		\item Hàm số $y=f(x)$ nghịch biến trên các khoảng $(-\infty; 1)$ và $(3; +\infty)$.  
		\item Đồ thị hàm số có tiệm cận đứng là $x = 2$ và tiệm cận xiên là $y = x + 1$.  
		\item Đồ thị hàm số có điểm cực đại là $(1; 1)$ và điểm cực tiểu là $(3; 5)$.  
		\item Trên khoảng $(2; +\infty)$, giá trị nhỏ nhất của hàm số $y=f(x)$ bằng $5$.  
	\end{enumerate}
	
	\vspace{0.6cm}
	\noindent \textbf{Câu 37.} Cho hàm số $y = f(x) = \frac{-x^2+4x-4}{x-1}$ có đồ thị $(C)$ như hình vẽ dưới đây.
	\begin{center}
		\begin{tikzpicture}
			\begin{axis}[
				width=9cm, height=8cm,
				axis lines=center,
				xmin=-3, xmax=5, ymin=-2, ymax=7,
				xtick={1, 2}, ytick={3, 4},
				xlabel={$x$}, ylabel={$y$},
				every axis x label/.style={at={(ticklabel* cs:1)},anchor=west},
				every axis y label/.style={at={(ticklabel* cs:1)},anchor=south},
				restrict y to domain=-5:15,
				]
				\addplot [domain=-3:0.7, samples=50, blue, thick] {(-x^2+4*x-4)/(x-1)};
				\addplot [domain=1.3:5, samples=50, blue, thick] {(-x^2+4*x-4)/(x-1)};
				\draw[dashed, red, thick] (axis cs:1,-2) -- (axis cs:1,7);
				\draw[dashed, red, thick] (axis cs:-3,6) -- (axis cs:5,-2);
				\draw[dashed, black!60] (axis cs:0,0) -- (axis cs:0,4) -- (axis cs:0,4);
				\draw[dashed, black!60] (axis cs:2,0) -- (axis cs:2,0) -- (axis cs:0,0);
				\filldraw (axis cs:0,4) circle (1.5pt);
				\filldraw (axis cs:2,0) circle (1.5pt);
				\node[above, sloped] at (axis cs:-2, 3) {$y = -x + 3$};
			\end{axis}
		\end{tikzpicture}
	\end{center}
	Các mệnh đề sau đúng hay sai?
	\begin{enumerate}[label=\textbf{\alph*)}]
		\item Hàm số đồng biến trên các khoảng $(0; 1)$ và $(1; 2)$.  
		\item Đồ thị hàm số có tiệm cận đứng $x = 1$ và tiệm cận xiên $y = -x + 3$.  
		\item Điểm cực tiểu của đồ thị hàm số là $(0; 4)$.  
		\item Giá trị lớn nhất của hàm số trên khoảng $(1; +\infty)$ là $0$.  
	\end{enumerate}
	
	\vspace{0.6cm}
	\noindent \textbf{Câu 38.} Cho hàm số $y = f(x) = \frac{2x - 3}{x - 1}$ có đồ thị $(C)$ như hình vẽ dưới đây.
	\begin{center}
		\begin{tikzpicture}
			\begin{axis}[
				width=9cm, height=8cm,
				axis lines=center,
				xmin=-4, xmax=5, ymin=-4, ymax=6,
				xtick={1, 2, 3}, ytick={1, 2},
				xlabel={$x$}, ylabel={$y$},
				every axis x label/.style={at={(ticklabel* cs:1)},anchor=west},
				every axis y label/.style={at={(ticklabel* cs:1)},anchor=south},
				restrict y to domain=-10:10,
				]
				\addplot [domain=-4:0.7, samples=50, blue, thick] {(2*x-3)/(x-1)};
				\addplot [domain=1.3:5, samples=50, blue, thick] {(2*x-3)/(x-1)};
				\draw[dashed, red, thick] (axis cs:1,-4) -- (axis cs:1,6);
				\draw[dashed, red, thick] (axis cs:-4,2) -- (axis cs:5,2);
				\draw[dashed, black!60] (axis cs:2,0) -- (axis cs:2,1) -- (axis cs:0,1);
				\filldraw (axis cs:2,1) circle (1.5pt);
			\end{axis}
		\end{tikzpicture}
	\end{center}
	Các mệnh đề sau đúng hay sai?
	\begin{enumerate}[label=\textbf{\alph*)}]
		\item Hàm số đồng biến trên tập xác định $\mathbb{R} \setminus \{1\}$.  
		\item Đồ thị hàm số có tiệm cận ngang là $y = 2$ và tiệm cận đứng là $x = 1$.  
		\item Hàm số đã cho không có cực trị.  
		\item Giá trị lớn nhất của hàm số trên đoạn $[2; 3]$ là $3$.  
	\end{enumerate}
	
	\vspace{0.6cm}
	\noindent \textbf{Câu 39.} Cho hàm số $y = f(x) = \frac{1}{x^2 - 1}$ có đồ thị $(C)$ như hình vẽ dưới đây.
	\begin{center}
		\begin{tikzpicture}
			\begin{axis}[
				width=7cm, height=6.5cm,
				axis lines=center,
				xmin=-4, xmax=4, ymin=-3, ymax=3,
				xtick={-1, 1}, ytick={-1},
				xlabel={$x$}, ylabel={$y$},
				every axis x label/.style={at={(ticklabel* cs:1)},anchor=west},
				every axis y label/.style={at={(ticklabel* cs:1)},anchor=south},
				restrict y to domain=-10:10,
				]
				\addplot [domain=-4:-1.2, samples=50, blue, thick] {1/(x^2-1)};
				\addplot [domain=-0.8:0.8, samples=50, blue, thick] {1/(x^2-1)};
				\addplot [domain=1.2:4, samples=50, blue, thick] {1/(x^2-1)};
				\draw[dashed, red, thick] (axis cs:-1,-3) -- (axis cs:-1,3);
				\draw[dashed, red, thick] (axis cs:1,-3) -- (axis cs:1,3);
				\draw[dashed, red, thick] (axis cs:-4,0) -- (axis cs:4,0);
				\filldraw (axis cs:0,-1) circle (1.5pt);
			\end{axis}
		\end{tikzpicture}
	\end{center}
	Các mệnh đề sau đúng hay sai?
	\begin{enumerate}[label=\textbf{\alph*)}]
		\item Hàm số đạt cực đại tại $x = 0$ và giá trị cực đại của hàm số là $y = -1$.  
		\item Đồ thị hàm số có hai đường tiệm cận đứng là $x = 1, x = -1$ và một đường tiệm cận ngang là $y = 0$.  
		\item Giá trị nhỏ nhất của hàm số trên khoảng $(-1; 1)$ là $-1$.  
		\item Hàm số nghịch biến trên khoảng $(0; 1)$ và $(1; +\infty)$.  
	\end{enumerate}
	
	\vspace{0.6cm}
	\noindent \textbf{Câu 40.} Cho hàm số $y = f(x) = \frac{-x^2+2x-2}{x-1}$ có đồ thị $(C)$ như hình vẽ dưới đây.
	\begin{center}
		\begin{tikzpicture}
			\begin{axis}[
				width=9cm, height=8cm,
				axis lines=center,
				xmin=-3, xmax=5, ymin=-6, ymax=4,
				xtick={1, 2}, ytick={-2, 2},
				xlabel={$x$}, ylabel={$y$},
				every axis x label/.style={at={(ticklabel* cs:1)},anchor=west},
				every axis y label/.style={at={(ticklabel* cs:1)},anchor=south},
				restrict y to domain=-15:15,
				]
				\addplot [domain=-3:0.7, samples=50, blue, thick] {(-x^2+2*x-2)/(x-1)};
				\addplot [domain=1.3:5, samples=50, blue, thick] {(-x^2+2*x-2)/(x-1)};
				\draw[dashed, red, thick] (axis cs:1,-6) -- (axis cs:1,4);
				\draw[dashed, red, thick] (axis cs:-3,4) -- (axis cs:5,-4);
				\draw[dashed, black!60] (axis cs:2,0) -- (axis cs:2,-2) -- (axis cs:0,-2);
				\draw[dashed, black!60] (axis cs:0,2) -- (axis cs:0,2);
				\filldraw (axis cs:2,-2) circle (1.5pt);
				\filldraw (axis cs:0,2) circle (1.5pt);
				\node[above, sloped] at (axis cs:-1.5, 1) {$y = -x + 1$};
			\end{axis}
		\end{tikzpicture}
	\end{center}
	Các mệnh đề sau đúng hay sai?
	\begin{enumerate}[label=\textbf{\alph*)}]
		\item Đường tiệm cận xiên của đồ thị hàm số là đường thẳng $y = -x + 1$.  
		\item Đồ thị hàm số có điểm cực tiểu là $(0; 2)$ và điểm cực đại là $(2; -2)$.  
		\item Hàm số đồng biến trên các khoảng $(0; 1)$ và $(1; 2)$.  
		\item Trên đoạn $[2; 3]$, giá trị nhỏ nhất của hàm số bằng $-2$.  
	\end{enumerate}
	
	
	\vspace{0.4cm}
	\noindent \textbf{Câu 41.} Cho hàm số $y = \frac{x^2 - x + 3}{x + 1}$ có đồ thị là $(C)$. Các mệnh đề sau đúng hay sai?
	\begin{enumerate}[label=\textbf{\alph*)}]
		\item Đồ thị hàm số $(C)$ có hai đường tiệm cận gồm một đường đứng và một đường xiên.  
		\item Giao điểm của hai đường tiệm cận là điểm $I(-1; -3)$.  
		\item Khoảng cách từ gốc tọa độ $O$ đến đường tiệm cận xiên bằng $\sqrt{2}$.  
		\item Đường tiệm cận xiên của đồ thị hàm số đi qua điểm $M(0; -2)$.  
	\end{enumerate}
	
	\vspace{0.4cm}
	\noindent \textbf{Câu 42.} Cho hàm số $y = \frac{2x - 3}{x - 2}$ có đồ thị là $(C)$. Các mệnh đề sau đúng hay sai?
	\begin{enumerate}[label=\textbf{\alph*)}]
		\item Đồ thị hàm số $(C)$ có hai tiệm cận đứng.  
		\item Giao điểm của hai đường tiệm cận là điểm $I(2; 2)$.  
		\item Khoảng cách từ gốc tọa độ $O$ đến đường tiệm cận đứng bằng $2$.  
		\item Đường tiệm cận ngang của đồ thị hàm số đi qua điểm $A(1; 2)$.  
	\end{enumerate}
	
	\vspace{0.4cm}
	\noindent \textbf{Câu 43.} Cho hàm số $y = \frac{2x^2 + x - 1}{x + 1}$ có đồ thị là $(C)$. Các mệnh đề sau đúng hay sai?
	\begin{enumerate}[label=\textbf{\alph*)}]
		\item Đồ thị hàm số $(C)$ có tiệm cận đứng là đường thẳng $x = -1$.  
		\item Đồ thị hàm số $(C)$ không có bất kỳ đường tiệm cận nào.  
		\item Giới hạn tại điểm gián đoạn của hàm số là $\lim\limits_{x \to -1} y = -3$.  
		\item Đồ thị hàm số có một đường tiệm cận xiên là $y = 2x - 1$.  
	\end{enumerate}
	
	\vspace{0.4cm}
	\noindent \textbf{Câu 44.} Cho hàm số $y = \frac{2x^2 - 3x + 3}{x - 1}$ có đồ thị là $(C)$. Các mệnh đề sau đúng hay sai?
	\begin{enumerate}[label=\textbf{\alph*)}]
		\item Đồ thị hàm số $(C)$ có một tiệm cận đứng và một tiệm cận xiên.  
		\item Giao điểm của hai đường tiệm cận là điểm $I(1; 1)$.  
		\item Khoảng cách từ gốc tọa độ $O$ đến đường tiệm cận xiên bằng $\frac{\sqrt{5}}{5}$.  
		\item Đường tiệm cận xiên của đồ thị hàm số đi qua điểm $M(2; 3)$.  
	\end{enumerate}
	
	\vspace{0.4cm}
	\noindent \textbf{Câu 45.} Cho hàm số $y = \frac{3x^2 - 1}{x^2 - 4}$ có đồ thị là $(C)$. Các mệnh đề sau đúng hay sai?
	\begin{enumerate}[label=\textbf{\alph*)}]
		\item Đồ thị hàm số $(C)$ có tổng cộng ba đường tiệm cận.  
		\item Đồ thị hàm số $(C)$ có hai đường tiệm cận đứng và một đường tiệm cận ngang.  
		\item Khoảng cách giữa hai đường tiệm cận đứng bằng $4$.  
		\item Đường tiệm cận ngang của đồ thị hàm số đi qua điểm $M(0; 3)$.  
	\end{enumerate}
	
	\vspace{0.4cm}
	\noindent \textbf{Câu 46.} Cho hàm số $y = \frac{-x^2 + 3x - 1}{x - 2}$ có đồ thị là $(C)$. Các mệnh đề sau đúng hay sai?
	\begin{enumerate}[label=\textbf{\alph*)}]
		\item Đồ thị hàm số $(C)$ có hai đường tiệm cận gồm một đường đứng và một đường xiên.  
		\item Giao điểm của hai đường tiệm cận là điểm $I(2; -1)$.  
		\item Khoảng cách từ gốc tọa độ $O$ đến đường tiệm cận xiên bằng $\frac{\sqrt{2}}{2}$.  
		\item Đường tiệm cận xiên của đồ thị hàm số đi qua điểm $M(3; -2)$.  
	\end{enumerate}
	
	\vspace{0.4cm}
	\noindent \textbf{Câu 47.} Cho hàm số $y = \frac{\sqrt{x^2 + 1}}{x - 1}$ có đồ thị là $(C)$. Các mệnh đề sau đúng hay sai?
	\begin{enumerate}[label=\textbf{\alph*)}]
		\item Đồ thị hàm số $(C)$ có tổng cộng ba đường tiệm cận.  
		\item Giao điểm của tiệm cận đứng với các tiệm cận ngang lần lượt là $I_1(1; 1)$ và $I_2(1; -1)$.  
		\item Khoảng cách giữa hai đường tiệm cận ngang của đồ thị hàm số bằng $2$.  
		\item Đồ thị hàm số $(C)$ chỉ có duy nhất một đường tiệm cận ngang.  
	\end{enumerate}
	
	\vspace{0.4cm}
	\noindent \textbf{Câu 48.} Cho hàm số $y = \sqrt{x^2 + 4x + 5}$ có đồ thị là $(C)$. Các mệnh đề sau đúng hay sai?
	\begin{enumerate}[label=\textbf{\alph*)}]
		\item Đồ thị hàm số $(C)$ không có đường tiệm cận đứng và đường tiệm cận ngang.  
		\item Đồ thị hàm số $(C)$ có hai đường tiệm cận xiên là $y = x + 2$ (khi $x \to +\infty$) và $y = -x - 2$ (khi $x \to -\infty$).  
		\item Giao điểm của hai đường tiệm cận xiên của đồ thị hàm số nằm trên trục hoành.  
		\item Khoảng cách từ gốc tọa độ $O$ đến đường tiệm cận xiên $y = x + 2$ bằng $\sqrt{2}$.  
	\end{enumerate}
	
	
	\vspace{0.4cm}
	\noindent \textbf{Câu 41.} Cho hàm số $y = \frac{x^2 - x + 3}{x + 1}$ có đồ thị là $(C)$. Các mệnh đề sau đúng hay sai?
	\begin{enumerate}[label=\textbf{\alph*)}]
		\item Đồ thị hàm số $(C)$ có hai đường tiệm cận gồm một đường đứng và một đường xiên.  
		\item Giao điểm của hai đường tiệm cận là điểm $I(-1; -3)$.  
		\item Khoảng cách từ gốc tọa độ $O$ đến đường tiệm cận xiên bằng $\sqrt{2}$.  
		\item Đường tiệm cận xiên của đồ thị hàm số đi qua điểm $M(0; -2)$.  
	\end{enumerate}
	
	\vspace{0.4cm}
	\noindent \textbf{Câu 42.} Cho hàm số $y = \frac{2x - 3}{x - 2}$ có đồ thị là $(C)$. Các mệnh đề sau đúng hay sai?
	\begin{enumerate}[label=\textbf{\alph*)}]
		\item Đồ thị hàm số $(C)$ có hai tiệm cận đứng.  
		\item Giao điểm của hai đường tiệm cận là điểm $I(2; 2)$.  
		\item Khoảng cách từ gốc tọa độ $O$ đến đường tiệm cận đứng bằng $2$.  
		\item Đường tiệm cận ngang của đồ thị hàm số đi qua điểm $A(1; 2)$.  
	\end{enumerate}
	
	\vspace{0.4cm}
	\noindent \textbf{Câu 43.} Cho hàm số $y = \frac{2x^2 + x - 1}{x + 1}$ có đồ thị là $(C)$. Các mệnh đề sau đúng hay sai?
	\begin{enumerate}[label=\textbf{\alph*)}]
		\item Đồ thị hàm số $(C)$ có tiệm cận đứng là đường thẳng $x = -1$.  
		\item Đồ thị hàm số $(C)$ không có bất kỳ đường tiệm cận nào.  
		\item Giới hạn tại điểm gián đoạn của hàm số là $\lim\limits_{x \to -1} y = -3$.  
		\item Đồ thị hàm số có một đường tiệm cận xiên là $y = 2x - 1$.  
	\end{enumerate}
	
	\vspace{0.4cm}
	\noindent \textbf{Câu 44.} Cho hàm số $y = \frac{2x^2 - 3x + 3}{x - 1}$ có đồ thị là $(C)$. Các mệnh đề sau đúng hay sai?
	\begin{enumerate}[label=\textbf{\alph*)}]
		\item Đồ thị hàm số $(C)$ có một tiệm cận đứng và một tiệm cận xiên.  
		\item Giao điểm của hai đường tiệm cận là điểm $I(1; 1)$.  
		\item Khoảng cách từ gốc tọa độ $O$ đến đường tiệm cận xiên bằng $\frac{\sqrt{5}}{5}$.  
		\item Đường tiệm cận xiên của đồ thị hàm số đi qua điểm $M(2; 3)$.  
	\end{enumerate}
	
	\vspace{0.4cm}
	\noindent \textbf{Câu 45.} Cho hàm số $y = \frac{3x^2 - 1}{x^2 - 4}$ có đồ thị là $(C)$. Các mệnh đề sau đúng hay sai?
	\begin{enumerate}[label=\textbf{\alph*)}]
		\item Đồ thị hàm số $(C)$ có tổng cộng ba đường tiệm cận.  
		\item Đồ thị hàm số $(C)$ có hai đường tiệm cận đứng và một đường tiệm cận ngang.  
		\item Khoảng cách giữa hai đường tiệm cận đứng bằng $4$.  
		\item Đường tiệm cận ngang của đồ thị hàm số đi qua điểm $M(0; 3)$.  
	\end{enumerate}
	
	\vspace{0.4cm}
	\noindent \textbf{Câu 46.} Cho hàm số $y = \frac{-x^2 + 3x - 1}{x - 2}$ có đồ thị là $(C)$. Các mệnh đề sau đúng hay sai?
	\begin{enumerate}[label=\textbf{\alph*)}]
		\item Đồ thị hàm số $(C)$ có hai đường tiệm cận gồm một đường đứng và một đường xiên.  
		\item Giao điểm của hai đường tiệm cận là điểm $I(2; -1)$.  
		\item Khoảng cách từ gốc tọa độ $O$ đến đường tiệm cận xiên bằng $\frac{\sqrt{2}}{2}$.  
		\item Đường tiệm cận xiên của đồ thị hàm số đi qua điểm $M(3; -2)$.  
	\end{enumerate}
	
	\vspace{0.4cm}
	\noindent \textbf{Câu 47.} Cho hàm số $y = \frac{\sqrt{x^2 + 1}}{x - 1}$ có đồ thị là $(C)$. Các mệnh đề sau đúng hay sai?
	\begin{enumerate}[label=\textbf{\alph*)}]
		\item Đồ thị hàm số $(C)$ có tổng cộng ba đường tiệm cận.  
		\item Giao điểm của tiệm cận đứng với các tiệm cận ngang lần lượt là $I_1(1; 1)$ và $I_2(1; -1)$.  
		\item Khoảng cách giữa hai đường tiệm cận ngang của đồ thị hàm số bằng $2$.  
		\item Đồ thị hàm số $(C)$ chỉ có duy nhất một đường tiệm cận ngang.  
	\end{enumerate}
	
	\vspace{0.4cm}
	\noindent \textbf{Câu 48.} Cho hàm số $y = \sqrt{x^2 + 4x + 5}$ có đồ thị là $(C)$. Các mệnh đề sau đúng hay sai?
	\begin{enumerate}[label=\textbf{\alph*)}]
		\item Đồ thị hàm số $(C)$ không có đường tiệm cận đứng và đường tiệm cận ngang.  
		\item Đồ thị hàm số $(C)$ có hai đường tiệm cận xiên là $y = x + 2$ (khi $x \to +\infty$) và $y = -x - 2$ (khi $x \to -\infty$).  
		\item Giao điểm của hai đường tiệm cận xiên của đồ thị hàm số nằm trên trục hoành.  
		\item Khoảng cách từ gốc tọa độ $O$ đến đường tiệm cận xiên $y = x + 2$ bằng $\sqrt{2}$.  
	\end{enumerate}


	\vspace{0.4cm}
	\noindent \textbf{Câu 49.} Cho hàm số $y = f(x) = \frac{x^2 + 4x - 1}{x + 1}$ có đồ thị là $(C)$ và đường tiệm cận xiên là đường thẳng $\Delta: y = ax + b$.
	\begin{center}
		\begin{tikzpicture}
			\begin{axis}[
				width=9cm, height=8cm,
				axis lines=center,
				xmin=-6, xmax=4, ymin=-5, ymax=9,
				xtick={-3, -1}, ytick={2, 3, 5},
				xlabel={$x$}, ylabel={$y$},
				every axis x label/.style={at={(ticklabel* cs:1)},anchor=west},
				every axis y label/.style={at={(ticklabel* cs:1)},anchor=south},
				restrict y to domain=-15:15,
				]
				\addplot [domain=-6:-1.3, samples=50, blue, thick] {(x^2+4*x-1)/(x+1)};
				\addplot [domain=-0.7:4, samples=50, blue, thick] {(x^2+4*x-1)/(x+1)};
				\draw[dashed, red, thick] (axis cs:-1,-5) -- (axis cs:-1,9);
				\addplot [domain=-6:4, dashed, red, thick] {x + 3};
				\draw[dashed, black!60] (axis cs:-3,0) -- (axis cs:-3,0);
				\draw[dashed, black!60] (axis cs:-1,0) -- (axis cs:-1,2) -- (axis cs:0,2);
				\filldraw (axis cs:-1,2) circle (1.5pt);
				\filldraw (axis cs:-3,0) circle (1.5pt);
				\filldraw (axis cs:0,3) circle (1.5pt);
				\node[above, sloped] at (axis cs:2, 5.5) {$y = x + 3$};
			\end{axis}
		\end{tikzpicture}
	\end{center}
	Các mệnh đề sau đúng hay sai?
	\begin{enumerate}[label=\textbf{\alph*)}]
		\item Giao điểm của tiệm cận xiên $\Delta$ và trục hoành $Ox$ có hoành độ bằng $-3$.  
		\item Giao điểm của tiệm cận xiên $\Delta$ và tiệm cận đứng của đồ thị $(C)$ là điểm $I(-1; 2)$.  
		\item Diện tích tam giác $OAB$ tạo bởi tiệm cận xiên $\Delta$ với hai trục tọa độ bằng $4,5$.  
		\item Giá trị lớn nhất của hàm số của đường tiệm cận xiên trên đoạn $[0; 2]$ bằng $5$.  
	\end{enumerate}
	
	\vspace{0.4cm}
	\noindent \textbf{Câu 50.} Cho hàm số $y = f(x) = \frac{x^2 - 2x + 3}{x - 2}$ có đồ thị là $(C)$ và đường tiệm cận xiên là đường thẳng $\Delta: y = ax + b$.
	\begin{center}
		\begin{tikzpicture}
			\begin{axis}[
				width=9cm, height=8cm,
				axis lines=center,
				xmin=-2, xmax=6, ymin=-4, ymax=8,
				xtick={2}, ytick={2},
				xlabel={$x$}, ylabel={$y$},
				every axis x label/.style={at={(ticklabel* cs:1)},anchor=west},
				every axis y label/.style={at={(ticklabel* cs:1)},anchor=south},
				restrict y to domain=-15:15,
				]
				\addplot [domain=-2:1.5, samples=50, blue, thick] {(x^2-2*x+3)/(x-2)};
				\addplot [domain=2.5:6, samples=50, blue, thick] {(x^2-2*x+3)/(x-2)};
				\draw[dashed, red, thick] (axis cs:2,-4) -- (axis cs:2,8);
				\addplot [domain=-2:6, dashed, red, thick] {x};
				\draw[dashed, black!60] (axis cs:2,0) -- (axis cs:2,2) -- (axis cs:0,2);
				\filldraw (axis cs:2,2) circle (1.5pt);
				\node[above, sloped] at (axis cs:4, 4.5) {$y = x$};
			\end{axis}
		\end{tikzpicture}
	\end{center}
	Các mệnh đề sau đúng hay sai?
	\begin{enumerate}[label=\textbf{\alph*)}]
		\item Đồ thị hàm số $(C)$ không có tiệm cận ngang.  
		\item Đường tiệm cận xiên của đồ thị $(C)$ tạo với hai trục tọa độ một tam giác có diện tích bằng $2$.  
		\item Giao điểm hai tiệm cận của đồ thị $(C)$ nằm trên đường parabol $y = x^2 - 2$.  
		\item Đường tiệm cận xiên của đồ thị $(C)$ vuông góc với đường thẳng $x + y + 10 = 0$.  
	\end{enumerate}
	
	\vspace{0.4cm}
	\noindent \textbf{Câu 51.} Cho hàm số $y = f(x) = \frac{-2x^2 + 5x - 1}{x - 1}$ có đồ thị là $(C)$ và đường tiệm cận xiên là đường thẳng $\Delta: y = ax + b$.
	\begin{center}
		\begin{tikzpicture}
			\begin{axis}[
				width=9cm, height=8cm,
				axis lines=center,
				xmin=-3, xmax=5, ymin=-5, ymax=7,
				xtick={1, 1.5}, ytick={1, 3},
				xlabel={$x$}, ylabel={$y$},
				every axis x label/.style={at={(ticklabel* cs:1)},anchor=west},
				every axis y label/.style={at={(ticklabel* cs:1)},anchor=south},
				restrict y to domain=-15:15,
				]
				\addplot [domain=-3:0.7, samples=50, blue, thick] {(-2*x^2+5*x-1)/(x-1)};
				\addplot [domain=1.3:5, samples=50, blue, thick] {(-2*x^2+5*x-1)/(x-1)};
				\draw[dashed, red, thick] (axis cs:1,-5) -- (axis cs:1,7);
				\addplot [domain=-3:4, dashed, red, thick] {-2*x + 3};
				\draw[dashed, black!60] (axis cs:1,0) -- (axis cs:1,1) -- (axis cs:0,1);
				\filldraw (axis cs:1,1) circle (1.5pt);
				\filldraw (axis cs:1.5,0) circle (1.5pt);
				\filldraw (axis cs:0,3) circle (1.5pt);
				\node[above, sloped] at (axis cs:-0.5, 4.5) {$y = -2x + 3$};
			\end{axis}
		\end{tikzpicture}
	\end{center}
	Các mệnh đề sau đúng hay sai?
	\begin{enumerate}[label=\textbf{\alph*)}]
		\item Đường tiệm cận xiên của đồ thị $(C)$ là đường thẳng $\Delta: y = -2x + 3$.  
		\item Giao điểm hai đường tiệm cận của $(C)$ nằm trên đường thẳng phân giác $y = x$.  
		\item Diện tích tam giác $OAB$ tạo bởi tiệm cận xiên $\Delta$ với hai trục tọa độ bằng $\frac{9}{4}$.  
		\item Giá trị nhỏ nhất của hàm số của đường tiệm cận xiên trên đoạn $[1; 3]$ bằng $-3$.  
	\end{enumerate}
	
	\vspace{0.4cm}
	\noindent \textbf{Câu 52.} Cho hàm số $y = f(x) = \frac{3x^2 - 2x + 2}{x - 1}$ có đồ thị là $(C)$ và đường tiệm cận xiên là đường thẳng $\Delta: y = ax + b$.
	\begin{center}
		\begin{tikzpicture}
			\begin{axis}[
				width=9cm, height=8cm,
				axis lines=center,
				xmin=-2, xmax=5, ymin=-3, ymax=13,
				xtick={1}, ytick={1, 4},
				xlabel={$x$}, ylabel={$y$},
				every axis x label/.style={at={(ticklabel* cs:1)},anchor=west},
				every axis y label/.style={at={(ticklabel* cs:1)},anchor=south},
				restrict y to domain=-15:20,
				]
				\addplot [domain=-2:0.7, samples=50, blue, thick] {(3*x^2-2*x+2)/(x-1)};
				\addplot [domain=1.3:4, samples=50, blue, thick] {(3*x^2-2*x+2)/(x-1)};
				\draw[dashed, red, thick] (axis cs:1,-3) -- (axis cs:1,13);
				\addplot [domain=-2:4, dashed, red, thick] {3*x + 1};
				\draw[dashed, black!60] (axis cs:1,0) -- (axis cs:1,4) -- (axis cs:0,4);
				\filldraw (axis cs:1,4) circle (1.5pt);
				\filldraw (axis cs:0,1) circle (1.5pt);
				\node[above, sloped] at (axis cs:2.5, 9.5) {$y = 3x + 1$};
			\end{axis}
		\end{tikzpicture}
	\end{center}
	Các mệnh đề sau đúng hay sai?
	\begin{enumerate}[label=\textbf{\alph*)}]
		\item Đường tiệm cận xiên của đồ thị $(C)$ là đường thẳng $\Delta: y = 3x + 1$.  
		\item Tiệm cận xiên $\Delta$ tạo với hai trục tọa độ một tam giác có diện tích bằng $\frac{1}{6}$.  
		\item Giao điểm hai tiệm cận $I$ của đồ thị hàm số nằm trên parabol $y = x^2 + 3$.  
		\item Đường tiệm cận xiên $\Delta$ song song với đường thẳng có phương trình $y = 3x - 2026$.  
	\end{enumerate}
	
	\vspace{0.4cm}
	\noindent \textbf{Câu 53.} Cho hàm số $y = f(x) = \frac{x^2 + 1}{2x - 4}$ có đồ thị là $(C)$ và đường tiệm cận xiên là đường thẳng $\Delta: y = ax + b$.
	\begin{center}
		\begin{tikzpicture}
			\begin{axis}[
				width=9cm, height=8cm,
				axis lines=center,
				xmin=-4, xmax=6, ymin=-4, ymax=6,
				xtick={-2, 2}, ytick={1, 2},
				xlabel={$x$}, ylabel={$y$},
				every axis x label/.style={at={(ticklabel* cs:1)},anchor=west},
				every axis y label/.style={at={(ticklabel* cs:1)},anchor=south},
				restrict y to domain=-15:15,
				]
				\addplot [domain=-4:1.6, samples=50, blue, thick] {(x^2+1)/(2*x-4)};
				\addplot [domain=2.4:6, samples=50, blue, thick] {(x^2+1)/(2*x-4)};
				\draw[dashed, red, thick] (axis cs:2,-4) -- (axis cs:2,6);
				\addplot [domain=-4:6, dashed, red, thick] {0.5*x + 1};
				\draw[dashed, black!60] (axis cs:2,0) -- (axis cs:2,2) -- (axis cs:0,2);
				\filldraw (axis cs:2,2) circle (1.5pt);
				\filldraw (axis cs:-2,0) circle (1.5pt);
				\filldraw (axis cs:0,1) circle (1.5pt);
				\node[above, sloped] at (axis cs:4, 3.2) {$y = 0,5x + 1$};
			\end{axis}
		\end{tikzpicture}
	\end{center}
	Các mệnh đề sau đúng hay sai?
	\begin{enumerate}[label=\textbf{\alph*)}]
		\item Đường tiệm cận xiên của đồ thị là đường thẳng $\Delta: y = \frac{1}{2}x + 1$.  
		\item Giao điểm hai đường tiệm cận của đồ thị $(C)$ là điểm $I(2; 2)$.  
		\item Diện tích tam giác tạo bởi đường tiệm cận xiên $\Delta$ và hai trục tọa độ bằng $1$.  
		\item Đường tiệm cận xiên $\Delta$ vuông góc với đường thẳng $2x + y - 5 = 0$.  
	\end{enumerate}
	
	\vspace{0.4cm}
	\noindent \textbf{Câu 54.} Cho hàm số $y = f(x) = \sqrt{x^2 - 2x + 2}$ có đồ thị là $(C)$ và đường tiệm cận xiên bên phải là $\Delta_1: y = a_1x + b_1$, bên trái là $\Delta_2: y = a_2x + b_2$.
	\begin{center}
		\begin{tikzpicture}
			\begin{axis}[
				width=9cm, height=8cm,
				axis lines=center,
				xmin=-4, xmax=6, ymin=-2, ymax=6,
				xtick={1, 2}, ytick={1},
				xlabel={$x$}, ylabel={$y$},
				every axis x label/.style={at={(ticklabel* cs:1)},anchor=west},
				every axis y label/.style={at={(ticklabel* cs:1)},anchor=south},
				restrict y to domain=-5:15,
				]
				\addplot [domain=-4:6, samples=100, blue, thick] {sqrt(x^2 - 2*x + 2)};
				\addplot [domain=-1:6, dashed, red, thick] {x - 1};
				\addplot [domain=-4:3, dashed, red, thick] {-x + 1};
				\draw[dashed, black!60] (axis cs:1,0) -- (axis cs:1,0);
				\filldraw (axis cs:1,0) circle (1.5pt);
				\filldraw (axis cs:2,1) circle (1.5pt);
				\node[above, sloped] at (axis cs:4, 3.2) {$y = x - 1$};
				\node[above, sloped] at (axis cs:-2, 3.2) {$y = -x + 1$};
			\end{axis}
		\end{tikzpicture}
	\end{center}
	Các mệnh đề sau đúng hay sai?
	\begin{enumerate}[label=\textbf{\alph*)}]
		\item Đồ thị hàm số $(C)$ không có tiệm cận ngang và tiệm cận đứng.  
		\item Đồ thị hàm số $(C)$ có hai tiệm cận xiên vuông góc với nhau.  
		\item Giao điểm của hai đường tiệm cận xiên của đồ thị hàm số có tọa độ là $I(1; 0)$.  
		\item Đường tiệm cận xiên bên phải $\Delta_1$ của đồ thị hàm số đi qua điểm $M(2; 1)$.  
	\end{enumerate}
	
	\vspace{0.4cm}
	\noindent \textbf{Câu 55.} Cho hàm số $y = f(x) = \frac{-x^2 + x - 2}{x + 1}$ có đồ thị là $(C)$ và đường tiệm cận xiên là đường thẳng $\Delta: y = ax + b$.
	\begin{center}
		\begin{tikzpicture}
			\begin{axis}[
				width=9cm, height=8cm,
				axis lines=center,
				xmin=-6, xmax=4, ymin=-8, ymax=4,
				xtick={-1, 2}, ytick={2, 3},
				xlabel={$x$}, ylabel={$y$},
				every axis x label/.style={at={(ticklabel* cs:1)},anchor=west},
				every axis y label/.style={at={(ticklabel* cs:1)},anchor=south},
				restrict y to domain=-15:15,
				]
				\addplot [domain=-6:-1.3, samples=50, blue, thick] {(-x^2+x-2)/(x+1)};
				\addplot [domain=-0.7:4, samples=50, blue, thick] {(-x^2+x-2)/(x+1)};
				\draw[dashed, red, thick] (axis cs:-1,-8) -- (axis cs:-1,4);
				\addplot [domain=-6:4, dashed, red, thick] {-x + 2};
				\draw[dashed, black!60] (axis cs:-1,0) -- (axis cs:-1,3) -- (axis cs:0,3);
				\filldraw (axis cs:-1,3) circle (1.5pt);
				\filldraw (axis cs:2,0) circle (1.5pt);
				\filldraw (axis cs:0,2) circle (1.5pt);
				\node[above, sloped] at (axis cs:-4, 6.2) {$y = -x + 2$};
			\end{axis}
		\end{tikzpicture}
	\end{center}
	Các mệnh đề sau đúng hay sai?
	\begin{enumerate}[label=\textbf{\alph*)}]
		\item Đường tiệm cận xiên của đồ thị hàm số là đường thẳng $\Delta: y = -x + 2$.  
		\item Giao điểm hai đường tiệm cận của đồ thị $(C)$ nằm trên đường parabol $y = x^2 + 2$.  
		\item Diện tích tam giác tạo bởi đường tiệm cận xiên $\Delta$ và hai trục tọa độ bằng $2$.  
		\item Giá trị lớn nhất của hàm số của đường tiệm cận xiên trên đoạn $[-2; 0]$ bằng $2$.  
	\end{enumerate}
	
	\vspace{0.4cm}
	\noindent \textbf{Câu 56.} Cho hàm số $y = f(x) = \frac{2x^2 - x + 1}{x + 1}$ có đồ thị là $(C)$ và đường tiệm cận xiên là đường thẳng $\Delta: y = ax + b$.
	\begin{center}
		\begin{tikzpicture}
			\begin{axis}[
				width=9cm, height=8cm,
				axis lines=center,
				xmin=-6, xmax=4, ymin=-11, ymax=5,
				xtick={-1, 1.5}, ytick={-3, -5},
				xlabel={$x$}, ylabel={$y$},
				every axis x label/.style={at={(ticklabel* cs:1)},anchor=west},
				every axis y label/.style={at={(ticklabel* cs:1)},anchor=south},
				restrict y to domain=-20:10,
				]
				\addplot [domain=-6:-1.3, samples=50, blue, thick] {(2*x^2-x+1)/(x+1)};
				\addplot [domain=-0.7:4, samples=50, blue, thick] {(2*x^2-x+1)/(x+1)};
				\draw[dashed, red, thick] (axis cs:-1,-11) -- (axis cs:-1,5);
				\addplot [domain=-6:4, dashed, red, thick] {2*x - 3};
				\draw[dashed, black!60] (axis cs:-1,0) -- (axis cs:-1,-5) -- (axis cs:0,-5);
				\filldraw (axis cs:-1,-5) circle (1.5pt);
				\filldraw (axis cs:1.5,0) circle (1.5pt);
				\filldraw (axis cs:0,-3) circle (1.5pt);
				\node[above, sloped] at (axis cs:2, 1.2) {$y = 2x - 3$};
			\end{axis}
		\end{tikzpicture}
	\end{center}
	Các mệnh đề sau đúng hay sai?
	\begin{enumerate}[label=\textbf{\alph*)}]
		\item Đường tiệm cận xiên của đồ thị là đường thẳng $\Delta: y = 2x - 3$.  
		\item Giao điểm của hai đường tiệm cận nằm trên đường parabol $y = -5x^2$.  
		\item Diện tích tam giác tạo bởi đường tiệm cận xiên $\Delta$ và hai trục tọa độ bằng $2,25$.  
		\item Đường tiệm cận xiên $\Delta$ song song với đường thẳng có phương trình $y = 2x + 2026$.  
	\end{enumerate}
	
	\vspace{0.4cm}
	\noindent \textbf{Câu 57.} Cho hàm số $y = f(x) = \frac{x^2 - 4x + 5}{x - 2}$ có đồ thị là $(C)$ và đường tiệm cận xiên là đường thẳng $\Delta: y = ax + b$.
	\begin{center}
		\begin{tikzpicture}
			\begin{axis}[
				width=9cm, height=8cm,
				axis lines=center,
				xmin=-2, xmax=6, ymin=-5, ymax=5,
				xtick={2}, ytick={-2},
				xlabel={$x$}, ylabel={$y$},
				every axis x label/.style={at={(ticklabel* cs:1)},anchor=west},
				every axis y label/.style={at={(ticklabel* cs:1)},anchor=south},
				restrict y to domain=-15:15,
				]
				\addplot [domain=-2:1.6, samples=50, blue, thick] {(x^2-4*x+5)/(x-2)};
				\addplot [domain=2.4:6, samples=50, blue, thick] {(x^2-4*x+5)/(x-2)};
				\draw[dashed, red, thick] (axis cs:2,-5) -- (axis cs:2,5);
				\addplot [domain=-2:6, dashed, red, thick] {x - 2};
				\draw[dashed, black!60] (axis cs:2,0) -- (axis cs:2,0);
				\filldraw (axis cs:2,0) circle (1.5pt);
				\filldraw (axis cs:0,-2) circle (1.5pt);
				\node[above, sloped] at (axis cs:4, 2.2) {$y = x - 2$};
			\end{axis}
		\end{tikzpicture}
	\end{center}
	Các mệnh đề sau đúng hay sai?
	\begin{enumerate}[label=\textbf{\alph*)}]
		\item Đường tiệm cận đứng của đồ thị là đường thẳng $x = 2$ và tiệm cận xiên là $y = x - 2$.  
		\item Giao điểm của hai đường tiệm cận trùng với giao điểm của tiệm cận xiên và trục hoành $Ox$.  
		\item Diện tích tam giác tạo bởi đường tiệm cận xiên $\Delta$ và hai trục tọa độ bằng $2$.  
		\item Đường tiệm cận xiên $\Delta$ vuông góc với đường thẳng $y = -x + 5$.  
	\end{enumerate}
	
	\vspace{0.4cm}
	\noindent \textbf{Câu 58.} Cho hàm số $y = f(x) = \frac{-2x^2 + x + 1}{x - 2}$ có đồ thị là $(C)$ và đường tiệm cận xiên là đường thẳng $\Delta: y = ax + b$.
	\begin{center}
		\begin{tikzpicture}
			\begin{axis}[
				width=9cm, height=8cm,
				axis lines=center,
				xmin=-4, xmax=6, ymin=-15, ymax=3,
				xtick={2, -1.5}, ytick={-3, -7},
				xlabel={$x$}, ylabel={$y$},
				every axis x label/.style={at={(ticklabel* cs:1)},anchor=west},
				every axis y label/.style={at={(ticklabel* cs:1)},anchor=south},
				restrict y to domain=-25:10,
				]
				\addplot [domain=-4:1.6, samples=50, blue, thick] {(-2*x^2+x+1)/(x-2)};
				\addplot [domain=2.4:6, samples=50, blue, thick] {(-2*x^2+x+1)/(x-2)};
				\draw[dashed, red, thick] (axis cs:2,-15) -- (axis cs:2,3);
				\addplot [domain=-4:4, dashed, red, thick] {-2*x - 3};
				\draw[dashed, black!60] (axis cs:2,0) -- (axis cs:2,-7) -- (axis cs:0,-7);
				\filldraw (axis cs:2,-7) circle (1.5pt);
				\filldraw (axis cs:-1.5,0) circle (1.5pt);
				\filldraw (axis cs:0,-3) circle (1.5pt);
				\node[above, sloped] at (axis cs:-2, 1.2) {$y = -2x - 3$};
			\end{axis}
		\end{tikzpicture}
	\end{center}
	Các mệnh đề sau đúng hay sai?
	\begin{enumerate}[label=\textbf{\alph*)}]
		\item Đường tiệm cận xiên của đồ thị hàm số là $\Delta: y = -2x - 3$.  
		\item Giao điểm của hai đường tiệm cận nằm trên đường parabol $y = -x^2 - 3$.  
		\item Diện tích tam giác tạo bởi đường tiệm cận xiên $\Delta$ và hai trục tọa độ bằng $2,25$.  
		\item Giá trị lớn nhất của hàm số của đường tiệm cận xiên trên đoạn $[-1; 1]$ bằng $-1$.  
	\end{enumerate}
	
	
	\subsection{Câu hỏi Trả lời ngắn}
	
	\noindent \textbf{Câu 59.} Cho hàm số $y = f(x) = \frac{-x^2 + 4x - 2025}{x - 3}$ có đồ thị là $(C)$. Biết đồ thị hàm số có đường tiệm cận xiên dạng $y = ax + b$ và đường tiệm cận đứng dạng $x = c$. Tính giá trị của biểu thức $T = a + b + c$.
	\\[4pt] \textbf{Trả lời:} \otraloi
	
	\vspace{0.6cm}
	\noindent \textbf{Câu 60.} Biết đồ thị hàm số $y = \frac{x^2 - 3x + 5}{x - 4}$ có đường tiệm cận xiên là $y = ax + b$. Tính giá trị của biểu thức $S = a^2 + 3b$.
	\\[4pt] \textbf{Trả lời:} \otraloi
	
	\vspace{0.6cm}
	\noindent \textbf{Câu 61.} Biết đồ thị hàm số $y = \frac{2x^2 - 5x + 3}{x - 2}$ có đường tiệm cận xiên là $y = ax + b$. Tìm hệ số góc $a$ của đường tiệm cận xiên đó.
	\\[4pt] \textbf{Trả lời:} \otraloi
	
	\vspace{0.6cm}
	\noindent \textbf{Câu 62.} Gọi $I(x_0; y_0)$ là tọa độ giao điểm của đường tiệm cận đứng và tiệm cận xiên của đồ thị hàm số $y = 3x - 2 + \frac{2}{x - 1}$. Tính giá trị biểu thức $A = x_0 + 2y_0$.
	\\[4pt] \textbf{Trả lời:} \otraloi
	
	\vspace{0.6cm}
	\noindent \textbf{Câu 63.} Gọi $I(x_0; y_0)$ là tọa độ giao điểm của đường tiệm cận đứng và tiệm cận xiên của đồ thị hàm số $y = \frac{x^2 + 4x}{x - 1}$. Tìm giá trị của tung độ $y_0$.
	\\[4pt] \textbf{Trả lời:} \otraloi
	
	\vspace{0.6cm}
	\noindent \textbf{Câu 64.} Đồ thị hàm số $y = \frac{\sqrt{4x^2 + 1} - 2x}{x - 3}$ có các đường tiệm cận ngang có phương trình dạng $y = y_1$ và $y = y_2$. Tính tổng $y_1 + y_2$.
	\\[4pt] \textbf{Trả lời:} \otraloi
	
	\vspace{0.6cm}
	\noindent \textbf{Câu 65.} Đồ thị hàm số $y = \frac{x - 2}{|x^2 - 5x + 6|}$ có tất cả bao nhiêu đường tiệm cận đứng?
	\\[4pt] \textbf{Trả lời:} \otraloi
	
	\vspace{0.6cm}
	\noindent \textbf{Câu 66.} Tính tổng số đường tiệm cận đứng và tiệm cận ngang của đồ thị hàm số $y = \frac{2x^2 - 8}{x^2 - 3x + 2}$.
	\\[4pt] \textbf{Trả lời:} \otraloi
	
	\vspace{0.6cm}
	\noindent \textbf{Câu 67.} Gọi $I(x_0; y_0)$ là giao điểm của hai đường tiệm cận của đồ thị hàm số $y = \frac{3x - 1}{x + 2}$. Tính giá trị của biểu thức $S = x_0^2 + y_0^2$.
	\\[4pt] \textbf{Trả lời:} \otraloi
	
	\vspace{0.6cm}
	\noindent \textbf{Câu 68.} Gọi $I(x_0; y_0)$ là giao điểm của đường tiệm cận đứng và đường tiệm cận ngang của đồ thị hàm số $y = \frac{x^2 - x - 2}{x^2 - 1}$. Tính bình phương khoảng cách từ gốc tọa độ $O(0; 0)$ đến điểm $I$ (tức là tính $OI^2$).
	\\[4pt] \textbf{Trả lời:} \otraloi
	
	\vspace{0.6cm}
	\noindent \textbf{Câu 69.} Cho đồ thị hàm số $y = \frac{3x^2 + 5x - 1}{x + 2}$ có tiệm cận xiên là $\Delta: y = ax + b$. Tìm giá trị của hàm số $y = ax + b$ tại điểm $x = 5$.
	\\[4pt] \textbf{Trả lời:} \otraloi
	
	\vspace{0.6cm}
	\noindent \textbf{Câu 70.} Gọi $I(x_0; y_0)$ là tọa độ giao điểm của đường tiệm cận đứng và tiệm cận xiên của đồ thị hàm số $y = \frac{4x^2 - 10x + 7}{2x - 4}$. Tính tích giá trị $P = x_0 \cdot y_0$.
	\\[4pt] \textbf{Trả lời:} \otraloi
	
	\vspace{0.6cm}
	\noindent \textbf{Câu 71.} Đồ thị hàm số $y = \frac{x - 1}{x^2 - 3x + 2}$ có tổng cộng bao nhiêu đường tiệm cận đứng và tiệm cận ngang?
	\\[4pt] \textbf{Trả lời:} \otraloi
	
	\vspace{0.6cm}
	\noindent \textbf{Câu 72.} Tính bình phương khoảng cách từ điểm $M(1; 4)$ đến đường tiệm cận xiên của đồ thị hàm số $y = \frac{x^2 + 2x - 1}{x + 1}$.
	\\[4pt] \textbf{Trả lời:} \otraloi
	
	\vspace{0.6cm}
	\noindent \textbf{Câu 73.} Đồ thị hàm số $y = \sqrt{4x^2 + 4x + 2}$ có hai đường tiệm cận xiên dạng $\Delta_1: y = a_1x + b_1$ (khi $x \to +\infty$) và $\Delta_2: y = a_2x + b_2$ (khi $x \to -\infty$). Tính giá trị biểu thức $T = a_1 + b_1 + a_2 + b_2$.
	\\[4pt] \textbf{Trả lời:} \otraloi
	
	\vspace{0.6cm}
	\noindent \textbf{Câu 74.} Gọi $I(x_0; y_0)$ là giao điểm của hai đường tiệm cận của đồ thị hàm số $y = \frac{x^2 - x + 1}{x - 1}$. Tính giá trị biểu thức $S = x_0^2 + y_0^2$.
	\\[4pt] \textbf{Trả lời:} \otraloi
	
	
	
	% =====================
	% PHẦN 5: TỔNG KẾT
	% =====================
	\section{Tổng kết bài và bài tập về nhà}
	
	\begin{tcolorbox}[colback=orange!10!white, colframe=orange!80!black,
		fonttitle=\bfseries, title={Bảng tóm tắt công thức tìm nhanh Tiệm cận}]
		\renewcommand{\arraystretch}{1.3}
		\begin{tabular}{p{4cm} p{6cm} p{6cm}}
			\textbf{Loại tiệm cận} & \textbf{Định nghĩa giới hạn} & \textbf{Cách tìm nhanh} \\
			\hline
			\textbf{Tiệm cận Đứng (TCĐ)} & $\lim\limits_{x \to x_0^{\pm}} f(x) = \pm\infty$ & Nghiệm của mẫu nhưng không là nghiệm của tử. \\
			\textbf{Tiệm cận Ngang (TCN)} & $\lim\limits_{x \to \pm\infty} f(x) = y_0$ & So sánh bậc tử và mẫu khi $x \to \pm\infty$. \\
			\textbf{Tiệm cận Xiên (TCX)} & $\lim\limits_{x \to \pm\infty} [f(x) - (ax+b)] = 0$ & Chia đa thức lấy phần thương hoặc bấm máy tính hệ số $a, b$. \\
		\end{tabular}
	\end{tcolorbox}
	
	\vspace{3mm}
	\textbf{Bài tập về nhà:} Hoàn thành các câu hỏi trắc nghiệm của bài học trong sách bài tập và chuẩn bị trước bài "Khảo sát sự biến thiên và vẽ đồ thị hàm số".
	
\end{document}